%\newcommand{cosec~}{\operatorname{cosec}}
\chapter{Trigonometry}
\setcounter{figure}{1}
\setcounter{subfigure}{1}

Trigonometry deals with the relationship between the angles and sides
of a right-angled triangle. We will learn about trigonometric
functions, which form the basis of trigonometry.

\chapterstartvideo{VMbec}

\section{Trigonometry is useful}
\nopagebreak
There are many applications of trigonometry. Of particular value is
the technique of triangulation, which is used in astronomy to measure
the distances to nearby stars, in geography to measure distances
between landmarks, and in satellite navigation systems. GPS (the
global positioning system) would not be possible without
trigonometry. Other fields which make use of trigonometry include
acoustics, optics, analysis of financial markets, electronics,
probability theory, statistics, biology, medical imaging (CAT scans
and ultrasound), chemistry, cryptology, meteorology, oceanography,
land surveying, architecture, phonetics, engineering, computer
graphics and game development.

\section{Similarity of triangles}
\nopagebreak
If $\triangle ABC$ is similar to $\triangle DEF$, then this is written as:

\begin{equation*}
  \triangle ABC\ |||\ \triangle DEF
\end{equation*}

\setcounter{subfigure}{0}
\begin{figure}[H] % horizontal\label{m39405*id78218}% \label{m39405*secfhsst!!!underscore!!!id95}
%             \subsection{  Discussion : Uses of Trigonometry }
%             \nopagebreak
%             
%       \label{m39405*id78126}Select one of the fields that uses trigonometry from the list given above and write a 1-page report describing \textit{how} trigonometry is used in the field that you chose. \par 
\begin{center}
\begin{pspicture}(-1,-0.6)(6.6,3)
%\psgrid[gridcolor=gray]
\rput(1.2,0){\pstTriangle(0,0){A}(2; 25){B}(3; 125){C}}
\rput(5.2,0){\pstTriangle[unit=0.5](0,0){D}(2; 25){E}(3; 125){F}}
\end{pspicture}
\end{center}   
\end{figure}   

In similar triangles, it is possible to deduce ratios between corresponding sides:
\begin{align*}
  \frac{AB}{BC} &= \frac{DE}{EF} \\
  \frac{AB}{AC} &= \frac{DE}{DF} \\
  \frac{AC}{BC} &= \frac{DF}{EF} \\
  \frac{AB}{DE} &= \frac{BC}{EF} = \frac{AC}{DF}
\end{align*}

Another important fact about similar triangles $ABC$ and $DEF$ is that
the angle at vertex $A$ is equal to the angle at vertex $D$, the angle
at $B$ is equal to the angle at $E$, and the angle at $C$ is equal to
the angle at $F$.\par


\begin{align*}
  \hat A &= \hat D \\
  \hat B &= \hat E \\
  \hat C &= \hat F
\end{align*}

\begin{Investigation}{Ratios of similar triangles}
\nopagebreak
Draw three similar triangles of different sizes, with each triangle having interior angles equal to ${30}^{\circ }$, ${90}^{\circ }$ and ${60}^{\circ }$ as shown below. Measure angles and lengths very accurately in order to fill in the table (round answers to one decimal place):\par 

\setcounter{subfigure}{0}
\scalebox{1} % Change this value to rescale the drawing.
{
\begin{pspicture}(0,-2.3817186)(8.567187,2.3817186)
\psline[linewidth=0.04cm](0.301875,-1.9001563)(1.942836,-1.9001563)
\psline[linewidth=0.04cm](1.942836,-1.9001563)(1.942836,0.84772635)
\psline[linewidth=0.04cm](0.301875,-1.9001563)(1.942836,0.8604187)
\psline[linewidth=0.04cm](2.641875,-1.9201562)(4.641875,-1.9201562)
\psline[linewidth=0.04cm](4.641875,-1.9201562)(4.641875,1.5438437)
\psline[linewidth=0.04cm](2.641875,-1.9201562)(4.641875,1.5598438)
\psline[linewidth=0.04cm](5.521875,-1.9201562)(7.801521,-1.9201562)
\psline[linewidth=0.04cm](7.801521,-1.9201562)(7.801521,1.9617095)
\psline[linewidth=0.04cm](5.521875,-1.9201562)(7.801521,1.9796396)
\usefont{T1}{ptm}{m}{n}
\rput(0.26546875,-2.1101563){$C$}
\usefont{T1}{ptm}{m}{n}
\rput(2.8054688,-2.1901562){$F$}
\usefont{T1}{ptm}{m}{n}
\rput(5.6254687,-2.1901562){$K$}
\usefont{T1}{ptm}{m}{n}
\rput(2.0454688,-2.1301563){$B$}
\usefont{T1}{ptm}{m}{n}
\rput(4.7654686,-2.1901562){$E$}
\usefont{T1}{ptm}{m}{n}
\rput(8.105469,-2.2101562){$H$}
\usefont{T1}{ptm}{m}{n}
\rput(2.1254687,1.0898438){$A$}
\usefont{T1}{ptm}{m}{n}
\rput(4.865469,1.8098438){$D$}
\usefont{T1}{ptm}{m}{n}
\rput(8.105469,2.1698437){$G$}
\psframe[linewidth=0.04,dimen=outer](1.961875,-1.5601562)(1.601875,-1.9201562)
\psframe[linewidth=0.04,dimen=outer](4.661875,-1.5801562)(4.301875,-1.9401562)
\psframe[linewidth=0.04,dimen=outer](7.801875,-1.5801562)(7.441875,-1.9401562)
\usefont{T1}{ptm}{m}{n}
\rput(1.64,-0.25){$30^{\circ}$}
% \pscircle[linewidth=0.0060,dimen=outer](1.811875,0.02984375){0.03}
\usefont{T1}{ptm}{m}{n}
\rput(0.9,-1.6){$60^{\circ}$}
% \pscircle[linewidth=0.0060,dimen=outer](0.951875,-1.6101563){0.03}
\usefont{T1}{ptm}{m}{n}
\rput(3.3,-1.6){$60^{\circ}$}
% \pscircle[linewidth=0.0060,dimen=outer](3.291875,-1.6301563){0.03}
\usefont{T1}{ptm}{m}{n}
\rput(6.2,-1.6){$60^{\circ}$}
% \pscircle[linewidth=0.0060,dimen=outer](6.131875,-1.6301563){0.03}
\usefont{T1}{ptm}{m}{n}
\rput(4.3173437,0.4){$30^{\circ}$}
% \pscircle[linewidth=0.0060,dimen=outer](4.531875,0.64984375){0.03}
\usefont{T1}{ptm}{m}{n}
\rput(7.4973435,0.9){$30^{\circ}$}
% \pscircle[linewidth=0.0060,dimen=outer](7.711875,1.1498437){0.03}
\end{pspicture} 
}      
\par 
% \textbf{m39405*id78604}\par
\begin{table}[H]
% \begin{table}[H]
% \\ '' '0'
\begin{center}

\noindent
\setlength{\extrarowheight}{2pt}

\begin{tabular}{|m{2.5cm}|m{2.5cm}|m{2.5cm}|}\hline
% My position: 0
% my spanname: 
% my ct of spanspec: 0
% my column-count: 3
\multicolumn{3}{|c|}{Dividing lengths of sides (ratios)}
\\ \hline
%--------------------------------------------------------------------
\Large$\frac{AB}{BC}=$
&
\Large$\frac{AB}{AC}=$
&
\Large$\frac{CB}{AC}=$
% make-rowspan-placeholders
\\ \hline
%--------------------------------------------------------------------
\Large$\frac{DE}{EF}=$
&
\Large$\frac{DE}{DF}=$
&
\Large$\frac{FE}{DF}=$
% make-rowspan-placeholders
\\ \hline
%--------------------------------------------------------------------
\Large$\frac{GH}{HK}=$
&
\Large$\frac{GH}{GK}=$
&
\Large$\frac{KH}{GK}=$
% make-rowspan-placeholders
\\ \hline
%--------------------------------------------------------------------
\end{tabular}
\end{center}
% \begin{center}{\small\bfseries Table 14.1}\end{center}
% \begin{caption}{\small\bfseries Table 14.1}\end{caption}
\end{table}
What observations can you make about the ratios of the sides?
\end{Investigation}


    

\section{Defining the trigonometric ratios}
The ratios of similar triangles are used to define the trigonometric ratios.
Consider a right-angled triangle $ABC$.

\setcounter{subfigure}{0}
\begin{figure}[H] % horizontal\label{m39408*id79667}
\begin{center}
\begin{pspicture}(-1,-0.6)(4,3)
%\psgrid[gridcolor=gray]
\pstTriangle(0,0){A}(2.4;90){B}(3.35;0){C}
\pstRightAngle{B}{A}{C}
\pstMarkAngle{B}{C}{A}{$\theta$}
\pcline[linestyle=none](A)(B)
\aput{:U}{opposite}
\pcline[linestyle=none](B)(C)
\aput{:U}{hypotenuse}
\pcline[linestyle=none](A)(C)
\bput{:U}{adjacent}
\end{pspicture}
\end{center}
\end{figure}       
\par 

% \Note{In algebra, we often use the letter $x$ for our unknown variable (although we can use any other letter too, such as $a$, $b$, $k$, etc.). In trigonometry, we often use the Greek symbol $\theta $ for an unknown angle (we also use $\alpha $ , $\beta $ , $\gamma $ etc.).}


In the right-angled triangle, we refer to the lengths of the three
sides according to how they are placed in relation to the angle
$\theta $. The side opposite to the right-angle is labelled the
hypotenuse, the side opposite $\theta $ is labelled ``opposite'', the
side next to $\theta $ is labelled ``adjacent''. Note that the choice
of non-$90^{\circ}$ internal angle is arbitrary. You can choose either
internal angle and then define the adjacent and opposite sides
accordingly. However, the hypotenuse remains the same regardless of
which internal angle you are referring to (because it is {\em always}
opposite the right-angle and {\em always} the longest side).

We define the trigonometric ratios, sine, cosine and tangent of an
angle, as follows:

\begin{equation*}
\begin{array}{ccc}\hfill sin~ \theta & =& \frac{\mbox{opposite}}{\mbox{hypotenuse}}\hfill \\
\\
 \hfill cos~ \theta & =& \frac{\mbox{adjacent}}{\mbox{hypotenuse}}\hfill \\
\\
 \hfill tan~ \theta & =& \frac{\mbox{opposite}}{\mbox{adjacent}}\hfill 
\end{array}
\end{equation*}

These ratios, also known as trigonometric identities, relate the
lengths of the sides of a right-angled triangle to its interior
angles. These three ratios form the basis of trigonometry.

\par
\textbf{Important:} the definitions of opposite, adjacent and hypotenuse are only applicable when working with right-angled triangles! Always check to make sure your triangle has a right-angle before you use them, otherwise you will get the wrong answer.
\par
You may also hear people saying ``Soh Cah Toa''. This is a mnemonic
technique for remembering the trigonometric ratios:

% \textbf{m39408*id79953}\par
\begin{table}[H]
% \begin{table}[H]
% \\ '' '0'
\begin{center}

\noindent

\begin{tabular}{|l|}\hline

$\mbox{\textbf{s}in}~\theta=\frac{\mbox{\textbf{o}pposite}}{\mbox{\textbf{h}ypotenuse}} $
% make-rowspan-placeholders

\\  \hline
%--------------------------------------------------------------------

$\mbox{\textbf{c}os}~\theta=\frac{\mbox{\textbf{a}djacent}}{\mbox{\textbf{h}ypotenuse}} $
% make-rowspan-placeholders

\\ \hline
%--------------------------------------------------------------------

$\mbox{\textbf{t}an}~\theta=\frac{\mbox{\textbf{o}pposite}}{\mbox{\textbf{a}djacent}} $
% make-rowspan-placeholders

\\   \hline
%--------------------------------------------------------------------
\end{tabular}
\end{center}
% \begin{center}{\small\bfseries Table 14.2}\end{center}
% \begin{caption}{\small\bfseries Table 14.2}\end{caption}
\end{table}

\par
\section{Reciprocal functions}
Each of the three trigonometric functions has a reciprocal. The
reciprocals, cosecant, secant and cotangent, are defined as follows:
\begin{equation*}
\begin{array}{ccc}cosec~ \theta & =& \dfrac{1}{sin~ \theta } \vspace{3pt}\\
 sec~ \theta & =& \dfrac{1}{cos~ \theta } \vspace{3pt}\\
 cot~ \theta & =& \dfrac{1}{tan~ \theta }
\end{array}
\end{equation*}

We can also define these reciprocals for any right-angled triangle:
\begin{equation*}
\begin{array}{ccc}\hfill cosec~ \theta & =& \frac{\mbox{hypotenuse}}{\mbox{opposite}}\hfill \vspace{3pt}\\
 \hfill sec~ \theta & =& \frac{\mbox{hypotenuse}}{\mbox{adjacent}}\hfill \vspace{3pt}\\
 \hfill cot~ \theta & =& \frac{\mbox{adjacent}}{\mbox{opposite}}\hfill 
\end{array}
\end{equation*}

Note that:
\begin{equation*}
\begin{array}{ccc}sin~ \theta \times cosec~ \theta & =& 1 \vspace{3pt}\\
 cos~ \theta \times sec~ \theta & =& 1 \vspace{3pt}\\
tan~ \theta \times cot~ \theta & =& 1
\end{array}
\end{equation*}

\par
\mindsetvid{Discovering the tangent ratio}{VMber}
% 
% 
% \begin{activity}{Calculator work}
% 
% It is very important that you are familiar with your calculator's trigonometric buttons and how they operate.
% Use the following examples to make sure that you are using your calculator correctly:
% \begin{enumerate}[noitemsep, label=\textbf{\arabic*}. ] 
% \item $cos~ 30^{\circ} = 0,866$
% \item $sin~ 90^{\circ} = 1$
% \item $tan~ 60^{\circ} = 1,732$
% \item $4 sin~ 45^{\circ} = 2,828$
% \item $\frac{1}{3} cos~ 60^{\circ} = 0,67$
% \end{enumerate}
% 
% \end{activity}
\par
\textbf{Important:} most scientific calculators are quite similar but these steps might differ depending on the calculator you use. Make sure your calculator is in ``degrees'' mode.


\begin{wex}{Using your calculator}
{
\begin{minipage}{\textwidth}
Use your calculator to calculate the following (correct to $2$ decimal places):
  \begin{enumerate}[itemsep=4pt, label=\textbf{\arabic*}. ] 
  \item $cos~ 48^{\circ}$
  \item $2 ~sin~ 35^{\circ}$
  \item $tan^{2}~ 81^{\circ}$
  \item $3~ sin^{2}~ 72^{\circ}$
  \item $\frac{1}{4}~ cos~ 27^{\circ}$
  \item $\frac{5}{6}~ tan~ 34^{\circ}$
  \end{enumerate}
\end{minipage}
}
{
\westep{}
Press \fbox{cos} \fbox{48} \fbox{\LARGE =} $0,67$

\westep{}
Press \fbox{2} \fbox{sin} \fbox{35} \fbox{\LARGE =} $1,15$

\westep{}
Press \fbox{(} \fbox{tan} \fbox{81} \fbox{)} \fbox{x$^{2}$}  \fbox{\LARGE =} $39,86$
\\
OR\\
Press \fbox{tan} \fbox{81} \fbox{\LARGE =} \fbox{ANS} \fbox{x$^{2}$} \fbox{\LARGE =} $39,86$

\westep{}
Press \fbox{3} \fbox{(} \fbox{sin} \fbox{72} \fbox{)} \fbox{x$^{2}$} \fbox{\LARGE =} $2,71$
\\
OR\\
Press \fbox{sin} \fbox{72} \fbox{\LARGE =} \fbox{ANS} \fbox{x$^{2}$} \fbox{\LARGE =} \fbox{ANS} \fbox{$\times$} \fbox{3}

\westep{}
Press \fbox{(} \fbox{1} \fbox{$\div$} \fbox{4} \fbox{)} \fbox{cos} \fbox{27} \fbox{\LARGE =} $0,22$
\\
OR\\
Press \fbox{cos} \fbox{27} \fbox{\LARGE =} \fbox{ANS} \fbox{$\div$} \fbox{4} \fbox{\LARGE =} $0,22$

\westep{}
Press \fbox{(} \fbox{5} \fbox{$\div$} \fbox{6} \fbox{)} \fbox{tan} \fbox{34} \fbox{\LARGE =} $0,56$
\\
OR\\
Press \fbox{tan} \fbox{34} \fbox{\LARGE =} \fbox{ANS} \fbox{$\times$} \fbox{5} \fbox{$\div$} \fbox{6} \fbox{\LARGE =} $0,56$
}
\end{wex}

\begin{wex}
{Calculator work}
{If $x=25^{\circ}$ and $y=65^{\circ}$, use your calculator to determine whether the following statement is true or false:
\begin{equation*}
  sin~^{2} x + cos~^{2} (90^{\circ}-y) = 1
\end{equation*}
}
{
\westep{Calculate the left hand side of the equation}
Press \fbox{(} \fbox{sin} \fbox{25} \fbox{)} \fbox{x$^{2}$} \fbox{\LARGE +} \fbox{(} \fbox{cos} \fbox{(} \fbox{90} \fbox{\LARGE -} \fbox{65} \fbox{)} \fbox{)} \fbox{x$^{2}$}  \fbox{\LARGE =} $1$


\westep{Write the final answer}
LHS = RHS therefore the statement is true.
}
\end{wex}

\begin{exercises}{}{
\begin{enumerate}[noitemsep, label=\textbf{\arabic*}. ]
\item In each of the following triangles, state whether $a$, $b$ and $c$ are the hypotenuse, opposite or adjacent sides of the triangle with respect to $\theta$. 
% \setcounter{subfigure}{0}
\begin{center}
\scalebox{0.85} % Change this value to rescale the drawing.
{
\begin{pspicture}(0,-3.7754679)(20.177345,3.9279697)
\rput (1,2){\textbf{(a)}}
\rput (5.5,2){\textbf{(b)}}
\rput (10.5,2){\textbf{(c)}}
\rput (1.5,-1.5){\textbf{(d)}}
\rput (6,-1.5){\textbf{(e)}}
\rput (9.5,-1.5){\textbf{(f)}}
\psdots[dotsize=0.027999999](2.9173439,-1.5720303)
\psline[linewidth=0.04cm](5.517344,2.5079696)(7.9173436,0.107969664)
\psline[linewidth=0.04cm](7.9173436,0.107969664)(9.0173435,1.2079697)
\psline[linewidth=0.04cm](5.517344,2.5079696)(9.0173435,1.2079697)
\psline[linewidth=0.04cm](7.717344,0.30796966)(7.9173436,0.5079697)
\psline[linewidth=0.04cm](7.9173436,0.5079697)(8.117344,0.30796966)
\psline[linewidth=0.04cm](11.0573435,2.2679696)(11.0573435,0.067969665)
\psline[linewidth=0.04cm](11.0573435,0.067969665)(14.757344,0.067969665)
\psline[linewidth=0.04cm](11.0573435,2.2679696)(14.757344,0.067969665)
\psline[linewidth=0.04cm](1.2373447,-3.3120303)(3.9373438,-0.71203035)
\psline[linewidth=0.04cm](3.9373438,-0.71203035)(3.9373438,-3.3120303)
\psline[linewidth=0.04cm](3.9373438,-3.3120303)(1.2373447,-3.3120303)
\psline[linewidth=0.04cm](7.477345,-1.2120303)(5.7773438,-3.3120303)
\psline[linewidth=0.04cm](7.477345,-1.2120303)(8.9773445,-2.4120302)
\psline[linewidth=0.04cm](8.9773445,-2.4120302)(5.7773438,-3.3120303)
\psline[linewidth=0.04cm](7.257343,-1.5120304)(7.537344,-1.7320304)
\psline[linewidth=0.04cm](7.517344,-1.7120303)(7.717344,-1.4320303)
\psline[linewidth=0.04cm](11.097343,0.38796967)(11.397344,0.38796967)
\psline[linewidth=0.04cm](11.397344,0.38796967)(11.397344,0.08796967)
\psline[linewidth=0.04cm](3.9173427,-3.0320303)(3.6173437,-3.0320303)
\psline[linewidth=0.04cm](3.6173437,-3.0320303)(3.6173437,-3.3320303)
\psline[linewidth=0.04cm](10.0173435,-1.7320304)(11.417343,-3.2320304)
\psline[linewidth=0.04cm](11.417343,-3.2320304)(13.0173435,-1.7320304)
\psline[linewidth=0.04cm](10.0173435,-1.7320304)(13.0173435,-1.7320304)
\psline[linewidth=0.04cm](11.217343,-3.0320303)(11.417343,-2.8320303)
\psline[linewidth=0.04cm](11.417343,-2.8320303)(11.617344,-3.0320303)
\usefont{T1}{ptm}{m}{n}
\rput(8.662344,1.1179696){$\theta$}
\usefont{T1}{ptm}{m}{n}
\rput(13.842344,0.29796967){$\theta$}
\usefont{T1}{ptm}{m}{n}
\rput(3.682343,-1.3020303){$\theta$}
\usefont{T1}{ptm}{m}{n}
\rput(6.362344,-2.9020302){$\theta$}
\usefont{T1}{ptm}{m}{n}
\rput(10.502343,-1.9420303){$\theta$}
\psline[linewidth=0.04cm](2.2973437,2.2879696)(0.09734377,0.18796967)
\psline[linewidth=0.04cm](0.09734377,0.18796967)(4.4973435,0.18796967)
\psline[linewidth=0.04cm](2.2973437,2.2879696)(4.4973435,0.18796967)
\psline[linewidth=0.04cm](2.0973437,2.0879698)(2.2973437,1.8879696)
\psline[linewidth=0.04cm](2.2973437,1.8879696)(2.4973438,2.0879698)
\usefont{T1}{ptm}{m}{n}
\rput(0.54234374,0.37796965){$\theta$}
\usefont{T1}{ptm}{m}{n}
\rput(1.1764063,1.5979697){$a$}
% \rput (1, 2){1.}
\usefont{T1}{ptm}{m}{n}
\rput(2.1778126,-0.10203034){$b$}
\usefont{T1}{ptm}{m}{n}
\rput(3.759219,1.2979697){$c$}
\usefont{T1}{ptm}{m}{n}
\rput(6.596406,1.0179696){$a$}
\usefont{T1}{ptm}{m}{n}
\rput(7.479219,2.1179698){$c$}
\usefont{T1}{ptm}{m}{n}
\rput(8.797812,0.41796967){$b$}
\usefont{T1}{ptm}{m}{n}
\rput(13.076406,1.4979696){$a$}
\usefont{T1}{ptm}{m}{n}
\rput(10.777813,1.1979697){$b$}
\usefont{T1}{ptm}{m}{n}
\rput(12.659219,-0.20203033){$c$}
\usefont{T1}{ptm}{m}{n}
\rput(2.896407,-3.6220303){$a$}
\usefont{T1}{ptm}{m}{n}
\rput(2.2978134,-1.8220303){$b$}
\usefont{T1}{ptm}{m}{n}
\rput(4.2792187,-2.3220303){$c$}
\usefont{T1}{ptm}{m}{n}
\rput(6.596406,-1.8220303){$a$}
\usefont{T1}{ptm}{m}{n}
\rput(7.7978134,-3.1220303){$b$}
\usefont{T1}{ptm}{m}{n}
\rput(8.579218,-1.7220303){$c$}
\usefont{T1}{ptm}{m}{n}
\rput(10.496406,-2.6220303){$a$}
\usefont{T1}{ptm}{m}{n}
\rput(12.397812,-2.8220303){$b$}
\usefont{T1}{ptm}{m}{n}
\rput(11.579218,-1.4220303){$c$}
\rput{-66.69996}(-0.02258055,0.761629){\psarc[linewidth=0.04](0.5673438,0.39796966){0.35}{25.785578}{138.26501}}
\rput{106.83747}(12.386083,-6.856496){\psarc[linewidth=0.04](8.737344,1.1679693){0.34}{32.561234}{147.21713}}
\rput{93.56388}(15.110441,-13.562675){\psarc[linewidth=0.04](13.927344,0.31796986){0.39}{24.333387}{127.3505}}
\psarc[linewidth=0.04](14.697344,3.1279697){0.0}{0.0}{180.0}
\psarc[linewidth=0.04](14.837344,3.2279696){0.0}{0.0}{180.0}
\rput{153.95303}(6.6359906,-3.7989676){\psarc[linewidth=0.04](3.757346,-1.1320316){0.5}{50.00946}{137.84839}}
\psarc[linewidth=0.04](20.097343,3.9079697){0.0}{0.0}{180.0}
\psarc[linewidth=0.04](20.117344,3.8879697){0.0}{0.0}{180.0}
\psarc[linewidth=0.04](20.157345,3.8279696){0.0}{0.0}{180.0}
\rput{-46.143967}(4.180272,3.6097476){\psarc[linewidth=0.04](6.3273435,-3.1020286){0.57}{58.40874}{128.71255}}
\rput{-116.039185}(16.474606,6.6625834){\psarc[linewidth=0.04](10.317342,-1.8120308){0.52}{49.23777}{123.63383}}
\end{pspicture}
}
\end{center}


\item Use your calculator to determine the value of the following (correct to $2$ decimal places):
\begin{enumerate}[noitemsep, label=\textbf{(\alph*)} ]
\begin{multicols}{2} 
% \setcounter{enumi}{6} 
\item $tan~ 65^{\circ}$
\item $sin~ 38^{\circ}$
\item $cos~ 74^{\circ}$
\item $sin~ 12^{\circ}$
\item $cos~ 26^{\circ}$
\item $tan~ 49^{\circ}$
\item $\frac{1}{4}~ cos~ 20^{\circ}$
\item $3 ~tan ~40^{\circ}$
\item $\frac{2}{3}~ sin~ 90^{\circ}$
\end{multicols}
\end{enumerate}

\item If $x=39^{\circ}$ and $y=21^{\circ}$, use a calculator to determine whether the following statements are true or false:
\begin{enumerate}[itemsep=4pt, label=\textbf{(\alph*)} ]
\begin{multicols}{2} 
% \setcounter{enumi}{15} 
\item $cos ~x + 2~ cos~ x = 3 ~cos~ x$
\item $cos ~2y = cos~ y + cos~ y$
\item $tan~ x = \dfrac{sin~ x}{cos~ x}$
\item $cos~ (x+y) = cos~ x + cos ~y$
\end{multicols}
\end{enumerate}


\item Complete each of the following (the first example has been done
  for you):
\begin{center}
% \begin{minipage}{0.25 \textwidth}
\setcounter{subfigure}{0}
\scalebox{1}{
\begin{pspicture}(0,-2.0990624)(4.21625,2.0990624)
\psline[linewidth=0.04cm](0.4775,-1.5590625)(3.7775,1.5409375)
\psline[linewidth=0.04cm](3.7775,1.5409375)(3.7775,-1.5590625)
\psline[linewidth=0.04cm](0.4775,-1.5590625)(3.7775,-1.5590625)
\psline[linewidth=0.04cm](3.4775,-1.5590625)(3.4775,-1.2590625)
\psline[linewidth=0.04cm](3.4775,-1.2590625)(3.7775,-1.2590625)
\rput(4.025,1.8759375){$A$}
\rput(3.874375,-1.9240625){ $B$}
\rput(0.10625,-1.8240625){$C$}
\end{pspicture}
}
% \end{minipage}
\end{center}

\begin{enumerate}[noitemsep, label=\textbf{(\alph*)} ]
\begin{multicols}{2}
% \setcounter{enumi}{19}
\item $sin~ \hat{A} = \frac{\mbox{\footnotesize opposite}}{\mbox{\footnotesize hypotenuse}}=\frac{CB}{AC}$
\item $cos~ \hat{A} = $
\item $tan~ \hat{A} = $
\item $sin~ \hat{C} = $
\item $cos~ \hat{C} = $
\item $tan~ \hat{C} = $
\end{multicols}
\end{enumerate}

\item Use the triangle below to complete the following:
\begin{center}
% \begin{minipage}{0.25\textwidth}
\scalebox{1} % Change this value to rescale the drawing.
{
\begin{pspicture}(0,-1.9864062)(4.449063,1.9664062)
\psline[linewidth=0.04cm](0.66,-1.4935937)(3.1364062,-1.4935937)
\psline[linewidth=0.04cm](3.1364062,-1.4935937)(3.1,1.9464062)
\psline[linewidth=0.04cm](3.08,1.9464062)(0.68,-1.4935937)
\psline[linewidth=0.04cm](2.7364063,-1.4935937)(2.7364063,-1.0935937)
\psline[linewidth=0.04cm](2.7364063,-1.0935937)(3.1364062,-1.0935937)
\usefont{T1}{ptm}{m}{n}
\rput(1.4889063,0.49640626){$2$}
\usefont{T1}{ptm}{m}{n}
\rput(2.0890625,-1.7835938){$1$}
\usefont{T1}{ptm}{m}{n}
\rput(1.3,-1.2){$60^{\circ}$}
\usefont{T1}{ptm}{m}{n}
\rput(2.85,1.1){$30^{\circ}$}
\usefont{T1}{ptm}{m}{n}
\rput(3.4545314,0.01640625){$\sqrt{3}$}
\end{pspicture}
}

% \end{minipage}
\end{center}
\\

\begin{enumerate}[noitemsep, label=\textbf{(\alph*)} ]
\begin{multicols}{2}
% \setcounter{enumi}{25}
\item $sin ~60^{\circ} = $
\item $cos ~60^{\circ} = $
\item $tan ~60^{\circ}= $
\item $sin ~30^{\circ}= $
\item $cos ~30^{\circ}= $
\item $tan~ 30^{\circ}= $
\end{multicols}
\end{enumerate}


\item Use the triangle below to complete the following:
\begin{center}
% \begin{minipage}{0.25\textwidth}
\scalebox{1} % Change this value to rescale the drawing.
{
\begin{pspicture}(0,-2.24)(5.0271873,2.22)
\psline[linewidth=0.04cm](0.361875,-1.7)(4.061875,-1.7)
\psline[linewidth=0.04cm](4.061875,-1.7)(4.061875,2.2)
\psline[linewidth=0.04cm](4.061875,2.2)(0.361875,-1.7)
\psline[linewidth=0.04cm](3.661875,-1.7)(3.661875,-1.3)
\psline[linewidth=0.04cm](3.661875,-1.3)(4.061875,-1.3)
\rput(2.2314062,-2.09){$1$}
\rput(1.1554687,-1.39){$45^{\circ}$}
\rput(3.7554688,1.41){$45^{\circ}$}
\rput(1.6054688,0.51){$\sqrt{2}$}
\rput(4.331406,0.21){$1$}
\end{pspicture} 
}
% \end{minipage}
\end{center}

\begin{enumerate}[noitemsep, label=\textbf{(\alph*)} ]

% \setcounter{enumi}{31}
\item $sin ~45^{\circ} = $
\item $cos~ 45^{\circ} = $
\item $tan~ 45^{\circ}= $

\end{enumerate}
\end{enumerate}
% Automatically inserted shortcodes - number to insert 6
\par \practiceinfo
\par \begin{tabular}[h]{cccccc}
% Question 1
(1.)	00ke	&
% Question 2
(2.)	00kf	&
% Question 3
(3.)	00kg	&
% Question 4
(4.)	00kh	&
% Question 5
(5.)	00ki	&
% Question 6
(6.)	00kj	\\ % End row of shortcodes
\end{tabular}
}
% Automatically inserted shortcodes - number inserted 6
\end{exercises}

\section{Special angles}
For most angles we need a calculator to calculate the values of
$sin~$, $cos~$ and $tan~$. However, we saw in the previous exercise
that we could work out these values for some special angles. The
values of the trigonometric functions for these special angles are
listed in the table below.

Remember that the lengths of the sides of a right-angled triangle must
obey the Theorem of Pythagoras: the square of the hypotenuse equals
the sum of the squares of the two other sides.

% \textbf{m39408*id80733}\par
\begin{table}[H]
% \begin{table}[H]
% \\ '' '0'
\begin{center}
\setlength{\extrarowheight}{2.5pt}

\begin{tabular}{|c|c|c|c|}\hline
$\theta$
&
${30}^{\circ}$
&
${45}^{\circ}$
&
${60}^{\circ}$


\\ \hline
%--------------------------------------------------------------------
$cos~ \theta$
 &
$\frac{\sqrt{3}}{2}$
&
$\frac{1}{\sqrt{2}}$
&
$\frac{1}{2}$

\\ \hline
%--------------------------------------------------------------------
$sin~ \theta $
&

$\frac{1}{2} $
&
$\frac{1}{\sqrt{2}}$
&
$\frac{\sqrt{3}}{2}$

%  make-rowspan-placeholders
\\ \hline
%--------------------------------------------------------------------
$tan~ \theta $
&

$\frac{1}{\sqrt{3}}$
&
$1$ &
$\sqrt{3}$



% make-rowspan-placeholders
\\ \hline
%--------------------------------------------------------------------
\end{tabular}
\end{center}
% \begin{center}{\small\bfseries Table 14.3}\end{center}
% \begin{caption}{\small\bfseries Table 14.3}\end{caption}
\end{table}


\begin{minipage}{0.5\textwidth}
\begin{center}
\scalebox{1} % Change this value to rescale the drawing.
{
\begin{pspicture}(0,-1.9864062)(4.449063,1.9664062)
\psline[linewidth=0.04cm](0.66,-1.4935937)(3.1364062,-1.4935937)
\psline[linewidth=0.04cm](3.1364062,-1.4935937)(3.1,1.9464062)
\psline[linewidth=0.04cm](3.08,1.9464062)(0.68,-1.4935937)
\psline[linewidth=0.04cm](2.7364063,-1.4935937)(2.7364063,-1.0935937)
\psline[linewidth=0.04cm](2.7364063,-1.0935937)(3.1364062,-1.0935937)
\usefont{T1}{ptm}{m}{n}
\rput(1.4889063,0.49640626){$2$}
\usefont{T1}{ptm}{m}{n}
\rput(2.0890625,-1.7835938){$1$}
\usefont{T1}{ptm}{m}{n}
\rput(1.3,-1.2){$60^{\circ}$}
\usefont{T1}{ptm}{m}{n}
\rput(2.85,1.1){\small$30^{\circ}$}
\usefont{T1}{ptm}{m}{n}
\rput(3.4545314,0.01640625){$\sqrt{3}$}
\end{pspicture}
}
\end{center}
\end{minipage}


\begin{minipage}{0.5\textwidth}
\begin{center}
\scalebox{1} % Change this value to rescale the drawing.
{
\begin{pspicture}(0,-2.24)(5.0271873,2.22)
\psline[linewidth=0.04cm](0.361875,-1.7)(4.061875,-1.7)
\psline[linewidth=0.04cm](4.061875,-1.7)(4.061875,2.2)
\psline[linewidth=0.04cm](4.061875,2.2)(0.361875,-1.7)
\psline[linewidth=0.04cm](3.661875,-1.7)(3.661875,-1.3)
\psline[linewidth=0.04cm](3.661875,-1.3)(4.061875,-1.3)
\rput(2.2314062,-2.09){$1$}
\rput(1.1554687,-1.39){$45^{\circ}$}
\rput(3.7554688,1.41){$45^{\circ}$}
\rput(1.6054688,0.51){$\sqrt{2}$}
\rput(4.331406,0.21){$1$}
\end{pspicture}
}
\end{center}
\end{minipage}

\par
These values are useful when we need to solve a problem involving trigonometric functions without using a calculator.\par 
\begin{exercises}{}{
\begin{enumerate}[itemsep=6pt, label=\textbf{\arabic*}. ] 
\item Calculate the value of the following without using a calculator:
\begin{enumerate}[noitemsep, label=\textbf{(\alph*)} ]
\item $sin~ 45^{\circ} \times cos~ 45^{\circ}$
\item $cos~ 60^{\circ} + tan~ 45^{\circ}$
\item $sin~ 60^{\circ} - cos~ 60^{\circ}$
\end{enumerate}

\item Use the table to  show that:
\begin{enumerate}[itemsep=5pt, label=\textbf{(\alph*)} ]
\item $\dfrac{sin~ 60^{\circ}}{cos~ 60^{\circ}} = tan~ 60^{\circ} $
\item $sin^{2} 45^{\circ}+ cos^{2} 45^{\circ} =1 $
\item $cos~ 30^{\circ} =\sqrt{1- sin^{2} 30^{\circ}}$
\end{enumerate}

\item Use the definitions of the trigonometric rations to answer the following questions:
\begin{enumerate}[noitemsep, label=\textbf{(\alph*)} ]
% \setcounter{enumi}{6}
\item Explain why $sin~ \alpha \leq 1$ for all values of $\alpha$.
\item Explain why $cos~ \beta$ has a maximum value of $1$.
\item Is there a maximum value for $tan~\gamma$?
\end{enumerate}
\end{enumerate}
% Automatically inserted shortcodes - number to insert 3
\par \practiceinfo
\par \begin{tabular}[h]{cccccc}
% Question 1
(1.)	00kk	&
% Question 2
(2.)	00km	&
% Question 3
(3.)	00kn	&
\end{tabular}
% Automatically inserted shortcodes - number inserted 3
}
\end{exercises}

\section{Solving trigonometric equations}
\mindsetvid{Choosing a ratio}{VMbfo}
\begin{wex}{Finding lengths}{Find the length of $x$ in the following right-angled triangle: \\
\begin{center}
\scalebox{1} 
{
\begin{pspicture}(0,-1.02)(4.655469,2.02)
\psline[linewidth=0.04](4.0,-1.0)(1.0,-1.0)(4.0,2.0)(4.0,-1.0)
\rput(2.0275,0.77){$100$}
\psarc[linewidth=0.04](1.0,-1.0){1.0}{0.0}{45.0}
\rput(1.6450001,-0.75){$50^{\circ}$}
\rput(4.355,0.41){$x$}
\psline[linewidth=0.04cm](3.6,-0.98)(3.6,-0.58)
\psline[linewidth=0.04cm](3.6,-0.58)(4.0,-0.58)
\end{pspicture} 
}
\end{center}
}
{
\westep{Identify the opposite and adjacent sides and the hypotenuse}
\begin{equation*}
\begin{array}{ccl}
 
\hfill sin~ \theta &=& \dfrac{\mbox{opposite}}{\mbox{hypotenuse}}  \hfill \vspace{5pt}\\
\hfill sin~  50^\circ &=& \dfrac{x}{100}  \hfill \\
\end{array}
\end{equation*}



\westep{Rearrange the equation to solve for $x$}
\begin{equation*}
 x=100 \times sin~ 50^{\circ}
\end{equation*}

\westep{Use your calculator to find the answer}
\begin{equation*}
x = 76,6
\end{equation*}
}
\end{wex}




% The following videos provide a summary of what you have learnt so far.
% \setcounter{subfigure}{0}
% \begin{figure}[H] % horizontal\label{m39408*trigonometry-1}
% \textnormal{Khan Academy video on trigonometry - 1}\vspace{.1in} \nopagebreak
% \label{m39408*yt-media}\label{m39408*yt-video}
% \raisebox{-5 pt}{ \includegraphics[width=0.5cm]{col11306.imgs/summary_www.png}} { (Video:  MG10102 )}
% \vspace{2pt}
% \vspace{.1in}
% \end{figure}       \par \label{m39408*eip-33}
% \setcounter{subfigure}{0}
% \begin{figure}[H] % horizontal\label{m39408*trigonometry-2}
% \textnormal{Khan Academy video on trigonometry - 2}\vspace{.1in} \nopagebreak
% \label{m39408*yt-media2}\label{m39408*yt-video2}
% \raisebox{-5 pt}{ \includegraphics[width=0.5cm]{col11306.imgs/summary_www.png}} { (Video:  MG10103 )}
% \vspace{2pt}
% \vspace{.1in}
% \end{figure}       


\begin{exercises}{}{
\begin{enumerate}[itemsep=5pt, label=\textbf{\arabic*}. ]
\item In each triangle find the length of the side marked with a letter. Give answers correct to $2$ decimal places.
  \begin{center}
    \scalebox{1}{% Change this value to rescale the drawing.
      \begin{pspicture}(0,-4.11)(11.1198435,5)
        \psline[linewidth=0.04](4.974369,2.612122)(2.6350498,3.9603012)(0.6377475,0.49464345)(4.974369,2.612122)
        \psline[linewidth=0.04cm](2.8949742,3.8105037)(2.7451768,3.550579)
        \psline[linewidth=0.04cm](2.7451768,3.550579)(2.4852521,3.7003772)
        \psline[linewidth=0.04](4.557088,-3.2752385)(4.589271,-0.5754303)(0.5895552,-0.52775127)(4.557088,-3.2752385)
        \psline[linewidth=0.04cm](4.5856953,-0.8754088)(4.2857165,-0.87183315)
        \psline[linewidth=0.04cm](4.2857165,-0.87183315)(4.2892923,-0.5718545)
        \psline[linewidth=0.04](7.0587263,3.3569472)(6.922255,0.6603981)(10.917143,0.45821914)(7.0587263,3.3569472)
        \psline[linewidth=0.04cm](6.9374194,0.9600146)(7.2370353,0.94485116)
        \psline[linewidth=0.04cm](7.2370353,0.94485116)(7.2218714,0.64523476)
        \psline[linewidth=0.04](7.9354,-3.9493167)(10.598035,-3.9603012)(10.613694,-0.16448934)(7.9354,-3.9493167)
        \psline[linewidth=0.04cm](10.340361,-3.959238)(10.341382,-3.7116852)
        \psline[linewidth=0.04cm](10.341382,-3.7116852)(10.599056,-3.7127483)
        % \usefont{T1}{ptm}{m}{n}
        \rput(0.12453125,3.7196066){\textbf{(a)}}
        % \usefont{T1}{ptm}{m}{n}
        \rput(6.133125,3.7396066){\textbf{(b)}}
        % \usefont{T1}{ptm}{m}{n}
        \rput(0.11421875,-0.26039344){\textbf{(c)}}
        % \usefont{T1}{ptm}{m}{n}
        \rput(6.1343746,-0.28039345){\textbf{(d)}}
        \usefont{T1}{ptm}{m}{n}
        \rput(4.197969,2.6346066){$37^\circ$}
        \usefont{T1}{ptm}{m}{n}
        \rput(3.0853126,1.2){$62$}
        \usefont{T1}{ptm}{m}{n}
        \rput(1.3824998,2.5846066){$a$}
        \usefont{T1}{ptm}{m}{n}
        \rput(9.804531,0.8146064){$23^\circ$}
        \usefont{T1}{ptm}{m}{n}
        \rput(8.595937,0.21460645){$21$}
        \usefont{T1}{ptm}{m}{n}
        \rput(6.663907,2.0646067){$b$}
        \usefont{T1}{ptm}{m}{n}
        \rput(4.2134376,-2.6453934){$ 55^\circ$}
        \usefont{T1}{ptm}{m}{n}
        \rput(2.166719,-2.1653934){$19$}
        \usefont{T1}{ptm}{m}{n}
        \rput(4.8759375,-1.8353934){$c$}
        \usefont{T1}{ptm}{m}{n}
        \rput(8.953907,-1.9653934){$33$}
        \usefont{T1}{ptm}{m}{n}
        \rput(10.22,-1.1453935){ $49^\circ$}
        \usefont{T1}{ptm}{m}{n}
        \rput(10.90422,-2.0153935){$d$}
      \end{pspicture}    
    }
  \end{center}

  \begin{center}
    \scalebox{1}{% Change this value to rescale the drawing.
      \begin{pspicture}(0,-3.8416457)(11.207031,3.8216705)
        \psline[linewidth=0.04](2.3197074,3.8016455)(1.0997375,1.3929795)(4.668132,-0.4143837)(2.3197074,3.8016455)
        \psline[linewidth=0.04cm](1.2352896,1.6606089)(1.5029193,1.5250567)
        \psline[linewidth=0.04cm](1.5029193,1.5250567)(1.3673671,1.2574271)
        % \usefont{T1}{ptm}{m}{n}
        \rput(0.11234375,3.5211668){\textbf{(e)}}
        % \usefont{T1}{ptm}{m}{n}
        \rput(6.084844,3.5211668){\textbf{(f)}}
        % \usefont{T1}{ptm}{m}{n}
        \rput(0.14375,-0.39883313){\textbf{(g)}}
        % \usefont{T1}{ptm}{m}{n}
        \rput(6.1403127,-0.498833){\textbf{(h)}}
        \psline[linewidth=0.04](1.0870312,-0.6888331)(1.0870312,-3.3888335)(5.0870314,-3.3888335)(1.0870312,-0.6888331)
        \psline[linewidth=0.04cm](1.0870312,-3.0888333)(1.3870313,-3.0888333)
        \psline[linewidth=0.04cm](1.3870313,-3.0888333)(1.3870313,-3.3888335)
        \psline[linewidth=0.04](10.705838,-3.6336784)(10.738021,-0.93386954)(6.738305,-0.88619083)(10.705838,-3.6336784)
        \psline[linewidth=0.04cm](10.734446,-1.2338486)(10.434467,-1.2302725)
        \psline[linewidth=0.04cm](10.434467,-1.2302725)(10.438043,-0.9302939)
        \psline[linewidth=0.04](11.200404,0.87534267)(9.325028,2.8177538)(6.4473815,0.039416883)(11.200404,0.87534267)
        \psline[linewidth=0.04cm](9.533403,2.6019304)(9.31758,2.393555)
        \psline[linewidth=0.04cm](9.31758,2.393555)(9.109204,2.6093786)
        \usefont{T1}{ptm}{m}{n}
        \rput(3.5887504,1.9261669){$e$}
        \usefont{T1}{ptm}{m}{n}
        \rput(3.8853126,0.3161668){$17^\circ$}
        \usefont{T1}{ptm}{m}{n}
        \rput(1.3671876,2.6761668){ $12$}
        \usefont{T1}{ptm}{m}{n}
        \rput(7.533282,0.51616687){$22^\circ$}
        \usefont{T1}{ptm}{m}{n}
        \rput(9.139219,0.2861668){$f$}
        \usefont{T1}{ptm}{m}{n}
        \rput(7.758125,1.7361668){ $31$}
        \usefont{T1}{ptm}{m}{n}
        \rput(3.028594,-1.683833){$32$}
        \usefont{T1}{ptm}{m}{n}
        \rput(4.1132817,-3.123833){$23^\circ$}
        \usefont{T1}{ptm}{m}{n}
        \rput(2.949844,-3.6138332){$g$}
        \usefont{T1}{ptm}{m}{n}
        \rput(10.928438,-2.2538335){$h$}
        \usefont{T1}{ptm}{m}{n}
        \rput(7.7667193,-1.2438331){ $30^\circ$}
        \usefont{T1}{ptm}{m}{n}
        \rput(8.245001,-2.3038335){$20$}
      \end{pspicture} 
    }
\end{center}  

\item Write down two ratios for each of the following in terms of the
  sides: $AB$; $BC$; $BD$; $AD$; $DC$ and $AC$:
  \begin{center}
    \scalebox{1}{% Change this value to rescale the drawing.
      \begin{pspicture}(0,-1.2515883)(4.02125,1.2515885)
        \psline[linewidth=0.04,fillstyle=solid](0.019002499,-1.2315885)(0.009700612,1.2315885)(4.00125,-1.1902716)(0.039492033,-1.2120903)(0.039492033,-1.2120903)(0.0,-1.1911051)
        \psline[linewidth=0.04](0.009700612,-0.9884115)(0.20970061,-0.9884115)(0.20970061,-1.2284116)
        \psline[linewidth=0.04,fillstyle=solid](1.3697007,0.4315885)(0.049700614,-1.1684115)(0.049700614,-1.1484115)
        \psline[linewidth=0.04](1.2297006,0.2715885)(1.0697006,0.3515885)(1.1897006,0.5115885)
        \usefont{T1}{ptm}{m}{n}
        \rput(0,-1.5){$A$}
        \rput(0, 1.5){$B$}
        \rput(1.4, 0.7){$C$}
        \rput(4, -1.5){$D$}
      \end{pspicture} 
    }
  \end{center}
  \begin{enumerate}[noitemsep, label=\textbf{(\alph*)} ]
  \item $sin~ \hat{B}$
  \item $cos~ \hat{D}$
  \item $tan~ \hat{B}$
  \end{enumerate}
  \vspace{10pt}
\item In $\triangle MNP$, $\hat{N}=90^{\circ}$, $MP=20$ and $\hat{P}=40^{\circ}$. Calculate $NP$ and $MN$ (correct to $2$ decimal places).
\end{enumerate}
\practiceinfo
\par 
\par \begin{tabular}[h]{cccccc}
(1.) 00kp& (2.) 00kq& (3.) 00mr \end{tabular}
}
\end{exercises}

\section{Finding an angle}
 \mindsetvid{What about the angle}{VMbfq}
\begin{wex}{Finding angles}
{Find the value of $\theta$ in the following right-angled triangle: \\
\begin{center}
\scalebox{1} % Change this value to rescale the drawing.
{
\begin{pspicture}(0,-0.9628125)(4.3034377,1.27)
\psline[linewidth=0.04](3.24,-0.57)(0.24,-0.57)(3.24,1.25)(3.24,-0.57)
\usefont{T1}{ptm}{m}{n}
\rput(1.9437499,-0.76){$100$}
\usefont{T1}{ptm}{m}{n}
\rput(3.8189065,0.28){$50$}
\usefont{T1}{ptm}{m}{n}
\rput(1.0326563,-0.3){$\theta$}
\psarc[linewidth=0.04](0.8,-0.45){0.8}{351.8699}{44.06081}
\psline[linewidth=0.04](3.26,-0.27)(2.92,-0.27)(2.92,-0.55)
\end{pspicture} 
}
\end{center}
}
{

\westep{Identify the opposite and adjacent sides and the hypotenuse}
In this case you have the opposite side and the hypotenuse for angle $\theta$. \\

\begin{equation*}
\begin{array}{ccl}
 
\hfill tan~ \theta &=& \dfrac{\mbox{opposite}}{\mbox{adjacent}}  \hfill \vspace{5pt}\\
\hfill tan~  \theta &=& \dfrac{50}{100} \hfill \\
\end{array}
\end{equation*}

\westep{Use your calculator to solve for $\theta$}
To solve for $\theta$, you will need to use the inverse function on your calculator: \vspace{10pt}
\\
Press \fbox{2ndF} \fbox{tan} \fbox{(} \fbox{50} \fbox{\div} \fbox{100} \fbox{)} \fbox{\LARGE =} $26,6$
\westep{Write the final answer}
\begin{equation*}
\theta = 26,6^{\circ}
\end{equation*}
}
\end{wex}

\begin{exercises}{}
{
   \begin{enumerate}[itemsep=5pt, label=\textbf{\arabic*}. ] 
\item Determine the angle (correct to $1$ decimal place):
    \begin{enumerate}[itemsep=3pt, label=\textbf{(\alph*)} ]
\begin{multicols}{2}
 \item $tan~ \theta = 1,7$
\item $sin~ \theta = 0,8$
\item $cos~ \alpha = 0,32$
\item $tan~ \theta = 5\frac{3}{4}$
\item $sin~ \theta = \frac{2}{3}$
\item $cos~ \gamma = 1,2$
\item $4 cos~ \theta = 3$
\item $cos~ 4\theta = 0,3$
\item $sin~ \beta + 2= 2,65$
\item $sin~ \theta = 0,8$
\item $3 tan~ \beta = 1$
\item $sin~ 3\alpha = 1,2$
\item $tan~ \frac{\theta}{3} = sin~ 48^{\circ}$
\item $\frac{1}{2} cos~ 2\beta = 0,3$
\item $2 sin~ 3\theta +1= 2,6$
\end{multicols}
\end{enumerate}

\item Determine $\alpha$ in the following right-angled triangles:
\begin{center}
\scalebox{1} % Change this value to rescale the drawing.
{
\begin{pspicture}(0,-4.032389)(7.3778124,4.17849)

\rput(0, 3.5){\textbf{(a)}}
\rput(3.5, 3.5){\textbf{(b)}}
\rput(0, 0.5){\textbf{(c)}}
\rput(3.5, 0.5){\textbf{(d)}}
\rput(0, -2.0){\textbf{(e)}}
\rput(3.5, -2.0){\textbf{(f)}}
\psline[linewidth=0.04,fillstyle=solid](4.128174,1.9359901)(4.129576,3.691884)(6.1603127,1.9559901)(4.142015,1.9498184)(4.142015,1.9498184)(4.115519,1.9649143)
\psline[linewidth=0.04,fillstyle=solid](0.5403125,3.7206802)(2.3103564,3.7223134)(0.5605087,1.1162983)(0.5542749,3.7072453)(0.5542749,3.7072453)(0.5694489,3.7330344)
\psline[linewidth=0.04](0.8203125,3.71599)(0.8203125,3.41599)(0.5403125,3.41599)
\psline[linewidth=0.04](4.132775,2.2306225)(4.429114,2.2301245)(4.4286456,1.9555349)
\psline[linewidth=0.04,fillstyle=solid](0.58817375,-1.2040099)(0.5895763,0.5518841)(2.6203125,-1.1840099)(0.60201496,-1.1901817)(0.60201496,-1.1901817)(0.57551914,-1.1750857)
\psline[linewidth=0.04](0.59277505,-0.90937746)(0.889114,-0.9098756)(0.88864577,-1.184465)
\psline[linewidth=0.04,fillstyle=solid](0.57534015,-1.9803641)(2.331222,-1.9737015)(0.6046739,-4.0123897)(0.5892318,-1.9941416)(0.5892318,-1.9941416)(0.6042059,-1.9675767)
\psline[linewidth=0.04](0.8699906,-1.983612)(0.8708536,-2.2799501)(0.5962649,-2.2807431)
\psline[linewidth=0.04,fillstyle=solid](6.6596694,0.2178756)(6.6803126,-0.6840099)(3.9803126,0.21599011)(6.645639,0.20423959)(6.645639,0.20423959)(6.6719236,0.18877967)
\psline[linewidth=0.04](6.6710052,-0.076665334)(6.3803124,-0.06400988)(6.3789563,0.20247632)
\psline[linewidth=0.04,fillstyle=solid](5.604427,-1.7804683)(6.7567835,-3.1053228)(4.0853443,-3.1304228)(5.603074,-1.7999867)(5.603074,-1.7999867)(5.632968,-1.7939644)
\psline[linewidth=0.04](5.794496,-2.005642)(5.570731,-2.199926)(5.3907113,-1.9925796)
\rput{-34.695152}(-0.9369147,0.7926677){\psarc[linewidth=0.024](0.8003125,1.8959901){0.3}{39.289406}{180.0}}
\usefont{T1}{ptm}{m}{n}
\rput(1.4359375,3.9229398){$4$}
\usefont{T1}{ptm}{m}{n}
\rput(0.29078126,2.4629397){$9$}
\usefont{T1}{ptm}{m}{n}
\rput(5.470625,3.02294){$13$}
\usefont{T1}{ptm}{m}{n}
\rput(5.034219,1.5829399){$7,5$}
\usefont{T1}{ptm}{m}{n}
\rput(1.75375,-0.057060193){$2,2$}
\usefont{T1}{ptm}{m}{n}
\rput(0.2096875,-0.4770602){$1,7$}
\usefont{T1}{ptm}{m}{n}
\rput(5.716875,0.4829398){$9,1$}
\usefont{T1}{ptm}{m}{n}
\rput(7.0603123,-0.3170602){$4,5$}
\usefont{T1}{ptm}{m}{n}
\rput(0.26703125,-2.8770602){$12$}
\usefont{T1}{ptm}{m}{n}
\rput(1.9092188,-3.1370602){$15$}
\usefont{T1}{ptm}{m}{n}
\rput(4.5784373,-2.3170602){$1$}
\usefont{T1}{ptm}{m}{n}
\rput(5.54375,-3.47706){$\sqrt{2}$}
\usefont{T1}{ptm}{m}{n}
\rput(6.2634373,-2.93706){$\alpha$}
\usefont{T1}{ptm}{m}{n}
\rput(0.7834375,-3.4970603){$\alpha$}
\usefont{T1}{ptm}{m}{n}
\rput(1.9234375,-0.9970602){$\alpha$}
\usefont{T1}{ptm}{m}{n}
\rput(5.1634374,-0.017060194){$\alpha$}
\usefont{T1}{ptm}{m}{n}
\rput(4.3034377,3.1429398){$\alpha$}
\usefont{T1}{ptm}{m}{n}
\rput(0.7634375,1.8029398){$\alpha$}
\rput{-162.53017}(7.4652467,7.66313){\psarc[linewidth=0.024](4.3213253,3.258065){0.34310362}{39.289406}{156.30495}}
\rput{-89.91282}(5.095528,5.0594177){\psarc[linewidth=0.024](5.081325,-0.02193504){0.34310362}{53.61859}{134.10115}}
\rput{-319.68652}(-0.17349519,-1.5915118){\psarc[linewidth=0.024](2.0811038,-1.0320795){0.42839816}{67.16096}{156.30495}}
\rput{-30.447884}(1.9160283,-0.06363264){\psarc[linewidth=0.024](0.8411037,-3.5520794){0.42839816}{67.16096}{156.30495}}
\rput{-319.68652}(-0.41087502,-4.8648257){\psarc[linewidth=0.024](6.4211035,-2.9920795){0.42839816}{67.16096}{156.30495}}
\end{pspicture} 

}
\end{center}
\end{enumerate}
% Automatically inserted shortcodes - number to insert 2
\par \practiceinfo
\par \begin{tabular}[h]{cccccc}
% Question 1
(1.)	00kr	&
% Question 2
(2.)	00ks	&
\end{tabular}
% Automatically inserted shortcodes - number inserted 2
}
\end{exercises}



\section{Two-dimensional problems}
Trigonometry was developed in ancient civilisations to solve practical problems such as building construction and navigating by the stars. We will show that trigonometry can also be used to solve some other practical problems. We use the trigonometric functions to solve problems in two dimensions that involve right-angled triangles. 
\par
\mindsetvid{Solving trig problems}{VMbkj}
\begin{wex}{Flying a kite}
{
\begin{minipage}{\textwidth}
Mandla flies a kite on a $17~$m string at an inclination of $63^{\circ}$.
\begin{enumerate}[noitemsep, label=\textbf{\arabic*}. ] 
 \item What is the height $h$ of the kite above the ground?
\item If Mandla's friend Sipho stands directly below the kite, calculate the distance $d$ between the two friends. 
\end{enumerate}
\end{minipage}
}
{
\westep{Make a sketch and identify the opposite and adjacent sides and the hypotenuse}
\begin{center}
\scalebox{1} % Change this value to rescale the drawing.
{
\begin{pspicture}(0,-1.7364062)(4.64375,1.7094063)
\psline[linewidth=0.025999999](4.0978127,1.6164063)(4.0978127,1.6964062)(4.0778127,-1.2435937)(0.4178126,-1.2435937)(4.0978127,1.6764063)
\usefont{T1}{ppl}{m}{n}
\rput(4.269219,0.1064063){$h$}
\usefont{T1}{ppl}{m}{n}
\rput(2.374844,-1.5335938){$d$}
\usefont{T1}{ppl}{m}{n}
\rput(1.1245313,-1.0135937){$63^{\circ}$}
\usefont{T1}{ppl}{m}{n}
\rput(2.31625,0.6464063){$17$}
\psline[linewidth=0.04](4.0778127,-0.98359364)(3.7978127,-0.98359364)(3.7978127,-1.2435937)
\end{pspicture} 
}
\end{center}

\westep{Use given information and appropriate ratio to solve for $h$ and $d$}
\begin{enumerate}[noitemsep, label=\textbf{\arabic*}. ] 
\item
\begin{equation*}
 \begin{array}{ccl}
\hfill sin~ 63^{\circ} &= & \frac{\mbox{\footnotesize opposite}}{\mbox{\footnotesize hypotenuse}} \hfill \vspace{5pt}\\
\hfill sin~ 63^{\circ} &=& \dfrac{h}{17} \hfill \\
\hfill \therefore h &= & 17 sin~ 63^{\circ} \hfill \\
\hfill &=& 15,15~\mbox{m} \hfill \\
   \end{array}
\end{equation*}

\item
\begin{equation*}
 \begin{array}{ccl}
\hfill cos~ 63^{\circ} &= & \frac{\mbox{\footnotesize adjacent}}{\mbox{\footnotesize hypotenuse}} \hfill \vspace{5pt}\\
\hfill cos~ 63^{\circ} &=& \dfrac{d}{17} \hfill \\
\hfill \therefore d &= & 17 cos~ 63^{\circ} \hfill \\
\hfill  &=& 7,72~\mbox{m} \hfill \\
\end{array}
\end{equation*}
\end{enumerate}
Note that the third side of the triangle can also be calculated using the Theorem of Pythagoras: $d^{2} = 17^{2} - h^{2}$.

\westep{Write final answers}
\begin{enumerate}[noitemsep, label=\textbf{\arabic*}. ] 
\item The kite is $15,15$ m above the ground.
\item Mandla and Sipho are $7,72$ m apart.
\end{enumerate}
}
\end{wex}

\begin{wex}
{Calculating angles}
{ABCD is a trapezium with $AB = 4$ cm, $CD = 6$ cm, $BC = 5$ cm and
  $AD = 5$ cm. Point $E$ on diagonal $AC$ divides the diagonal such
  that $AE = 3$ cm. $B\hat{E}C = 90^{\circ}$. Find $A\hat{B}C$.
}
{
\westep{Draw trapezium and label all given lengths on diagram. Indicate that $B\hat{E}C$ is a right-angle}
\begin{center}
\scalebox{1} % Change this value to rescale the drawing.
{
\begin{pspicture}(0,-2.8215597)(7.3365626,2.6459403)
\psline[linewidth=0.04](7.0446873,-2.3765597)(0.2846875,-2.3565598)(1.9846874,2.1834402)(5.5246873,2.1834402)(7.0246873,-2.3565598)
\psline[linewidth=0.04cm](2.0246875,2.1634402)(7.0246873,-2.3765597)
\psline[linewidth=0.04cm,linestyle=dashed,dash=0.16cm 0.16cm](5.5046873,2.1634402)(3.9646876,0.44344023)
\psline[linewidth=0.04](4.1846876,0.6834402)(4.4246874,0.48344022)(4.1646876,0.20344023)
\rput{124.910225}(10.204001,-1.1108128){\psarc[linewidth=0.04](5.3916802,2.1056097){0.48547706}{48.07423}{180.0}}
% \usefont{T1}{ptm}{m}{n}
\rput(5.7178125,2.4534402){$B$}
% \usefont{T1}{ptm}{m}{n}
\rput(7.23,-2.2465599){$C$}
% \usefont{T1}{ptm}{m}{n}
\rput(1.885,2.4534402){$A$}
% \usefont{T1}{ptm}{m}{n}
\rput(0.075,-2.2865598){$D$}
% \usefont{T1}{ptm}{m}{n}
\rput(3.7692187,0.15344024){$E$}
% \usefont{T1}{ptm}{m}{n}
\rput(3.7959375,2.4134402){$4$ cm}
% \usefont{T1}{ptm}{m}{n}
\rput(0.61078125,0.27344024){$5$ cm}
% \usefont{T1}{ptm}{m}{n}
\rput(6.790781,0.19344023){$5$ cm}
% \usefont{T1}{ptm}{m}{n}
\rput{-43.673504}(0.04067405,2.2905717){\rput(2.8509376,1.0734402){$3$ cm}}
% \usefont{T1}{ptm}{m}{n}
\rput(3.4921875,-2.6665597){$6$ cm}
\end{pspicture} 
}
\end{center}
     

\westep{Use $\triangle ABE$ and $\triangle CBE$ to determine the two angles at $\hat{B}$} 

\westep{Find the first angle, $A\hat{B}E$}
The hypotenuse and opposite side are given for both triangles,
therefore use the $sin~$ function: In $\triangle ABE$,

\begin{eqnarray*}
  sin~ A\hat{B}E & = & \frac{\mbox{\footnotesize opposite}}{\mbox{\footnotesize hypotenuse}} \\
                 & = & \dfrac{3}{4} \\
  \therefore A\hat{B}E & = & {48,6}^{\circ}
\end{eqnarray*}

\westep{Use Theorem of Pythagoras to determine $BE$}
In $\triangle ABE$, \\
\begin{eqnarray*}
  BE^{2} & = & AB^{2} - AE^{2} \\
        & = & 4^{2} - 3^{2} \\
        & = & 7 \\
  \therefore BE & = & \sqrt{7}
\end{eqnarray*}

\westep{Find the second angle $C\hat{B}E$}
In $\triangle CBE$ 
\begin{eqnarray*}
  cos~ C\hat{B}E & = & \frac{\mbox{\footnotesize adjacent}}{\mbox{\footnotesize hypotenuse}} \\
                 & = & \dfrac{\sqrt{7}}{5} \\
                 & = & 0,529 \\
  \therefore C\hat{B}E & = & {58,1}^{\circ}
\end{eqnarray*}

\westep{Calculate the sum of the angles}
\begin{equation*}
  A\hat{B}C = 48,6^{\circ} + 58,1^{\circ} = 106,7^{\circ}
\end{equation*}
}
\end{wex}

  
Another application is using trigonometry to find the height of a building. We could use a tape measure lowered from the roof, but this is impractical (and dangerous) for tall buildings. It is much more sensible to use trigonometry.\par 

\begin{wex}{Finding the height of a building}
{%
% \setcounter{subfigure}{0}
% \begin{figure}[htbp]
\begin{center}

\scalebox{1} % Change this value to rescale the drawing.
{
\begin{pspicture}(0,-2.4345315)(7.6,2.5717187)
\psframe[linewidth=0.028222222,dimen=outer,fillstyle=crosshatch,hatchwidth=0.028222222,hatchangle=180.0](7.6,2.0682812)(5.6,-1.9317188)
\psline[linewidth=0.028222222cm](0.6,-1.9317188)(5.6,-1.9317188)
\psline[linewidth=0.028222222cm,linestyle=dashed,dash=0.16cm 0.16cm](0.6,-1.9317188)(5.6,2.0682812)
\pswedge[linewidth=0.028222222](0.6,-1.9317188){0.6}{0.0}{38.66}
\usefont{T1}{ppl}{m}{n}
\rput(3.05,-2.231719){$100$ m}
\usefont{T1}{ppl}{m}{n}
\rput(1.6490624,-1.6317189){$38,7^\circ$}
\usefont{T1}{ppl}{m}{n}
\rput(0.33453125,-1.9317188){$Q$}
\usefont{T1}{ppl}{m}{n}
\rput(5.5490627,-2.231719){$B$}
\usefont{T1}{ppl}{m}{n}
\rput(5.5490627,2.3682814){$T$}
\usefont{T1}{ppl}{m}{n}
\rput(5.249062,0.06828115){$h$}
\psline[linewidth=0.04](5.61,-1.6082814)(5.25,-1.6082814)(5.25,-1.9282814)
\end{pspicture} 
}

\end{center}
\\

% \end{figure}
% \caption{Determining the height of a building using trigonometry.}
% \label{trig:height}

The given diagram shows a building of unknown height $h$. If we $(Q)$ walk $100$ m away from the building $(B)$ and measure the angle from the ground to the top of the building $(T)$, the angle is found to be $38,{7}^{\circ }$. This is called the angle of elevation.
We have a right-angled triangle and know the length of one side and an angle. We can therefore calculate the height of the building (correct to the nearest metre).}
{
\westep{Identify the opposite and adjacent sides and the hypotenuse}

\westep{}
In $\triangle QTB$,
\begin{equation*}
\begin{array}{ccl}\hfill tan~ 38,7^{\circ }& =& \frac{\mbox{\footnotesize opposite}}{\mbox{\footnotesize adjacent}}\hfill \vspace{4pt} \\
 & =& \dfrac{h}{100}\hfill
  \end{array}
\end{equation*}
\westep{Rearrange and solve for $h$}

\begin{equation*}
\begin{array}{ccl}

\hfill h& =& 100\times tan~ 38,7^{\circ }\hfill \\
& =& 80,1\hfill
  \end{array}
\end{equation*}

\westep{Write final answer}
The height of the building is $80$ m.
}
\end{wex}


\begin{wex}{Angles of elevation and depression}
{A block of flats is $200$ m away from a cellphone tower. Someone stands at $B$. They measure the angle from $B$ to the top of the tower $E$ to be $34^{\circ}$ (the angle of elevation). They then measure the angle from $B$ to the bottom of the tower $C$ to be $62^{\circ}$ (the angle of depression).
\\What is the height of the cellphone tower (correct to the nearest metre)?\\

\\
\begin{center}
\scalebox{1} % Change this value to rescale the drawing.
{
\begin{pspicture}(0,-3.893125)(10.969063,3.893125)
\definecolor{color194}{rgb}{0.2,0.2,0.2}
\definecolor{color353b}{rgb}{0.6,0.6,0.6}
\definecolor{color649b}{rgb}{0.8,0.8,1.0}
\psframe[linewidth=0.04,linecolor=color194,dimen=outer,fillstyle=solid](2.3,0.8496875)(0.18,-3.2103126)
\psframe[linewidth=0.04,linecolor=color194,dimen=outer,fillstyle=solid,fillcolor=color353b](0.62,-0.2903125)(0.4,-0.69031256)
\psframe[linewidth=0.04,linecolor=color194,dimen=outer,fillstyle=solid,fillcolor=color353b](1.32,-0.3103125)(1.1,-0.7103125)
\psframe[linewidth=0.04,linecolor=color194,dimen=outer,fillstyle=solid,fillcolor=color353b](0.62,-1.0303124)(0.4,-1.4303125)
\psframe[linewidth=0.04,linecolor=color194,dimen=outer,fillstyle=solid,fillcolor=color353b](0.62,0.4896875)(0.4,0.0896875)
\psframe[linewidth=0.04,linecolor=color194,dimen=outer,fillstyle=solid,fillcolor=color353b](1.32,0.4896875)(1.1,0.0896875)
\psframe[linewidth=0.04,linecolor=color194,dimen=outer,fillstyle=solid,fillcolor=color353b](2.02,0.4896875)(1.8,0.0896875)
\psframe[linewidth=0.04,linecolor=color194,dimen=outer,fillstyle=solid,fillcolor=color353b](2.02,-0.3103125)(1.8,-0.7103125)
\psframe[linewidth=0.04,linecolor=color194,dimen=outer,fillstyle=solid,fillcolor=color353b](1.32,-1.0303124)(1.1,-1.4303125)
\psframe[linewidth=0.04,linecolor=color194,dimen=outer,fillstyle=solid,fillcolor=color353b](2.02,-1.0303124)(1.8,-1.4303125)
\psframe[linewidth=0.04,linecolor=color194,dimen=outer,fillstyle=solid,fillcolor=color353b](0.62,-1.7303125)(0.4,-2.1303124)
\psframe[linewidth=0.04,linecolor=color194,dimen=outer,fillstyle=solid,fillcolor=color353b](1.32,-1.7303125)(1.1,-2.1303124)
\psframe[linewidth=0.04,linecolor=color194,dimen=outer,fillstyle=solid,fillcolor=color353b](2.02,-1.7303125)(1.8,-2.1303124)
\psframe[linewidth=0.04,linecolor=color194,dimen=outer,fillstyle=solid,fillcolor=color649b](0.5,-2.6103127)(0.2,-3.2103126)
\pscircle[linewidth=0.04,linecolor=color194,dimen=outer,fillstyle=gradient,gradlines=2000,gradmidpoint=1.0](0.35,-2.8603127){0.05}
\psframe[linewidth=0.04,linecolor=color194,dimen=outer,fillstyle=solid,fillcolor=color353b](10.2,2.2896874)(10.0,-3.2103126)
\psline[linewidth=0.04cm,linecolor=color194](10.1,3.2896874)(10.1,2.2896874)
\psline[linewidth=0.025999999cm,linecolor=color194](2.24,0.82968754)(9.98,-3.1703126)
\psline[linewidth=0.024cm,linecolor=color194](2.28,0.8496875)(10.12,3.2896874)
\psdots[dotsize=0.12,linecolor=color194](10.1,3.2896874)
\psline[linewidth=0.024cm,linecolor=color194,linestyle=dashed,dash=0.16cm 0.16cm](2.26,0.82968754)(10.02,0.82968754)
% \usefont{T1}{ptm}{m}{n}
\rput(5.776562,-3.6903126){$200$ m}
\usefont{T1}{ptm}{m}{n}
\rput(2.78,-2.7703125){$ A$}
\usefont{T1}{ptm}{m}{n}
\rput(2.2557812,1.1896875){$B$}
\usefont{T1}{ptm}{m}{n}
\rput(10.574532,-3.1703126){$C$}
\usefont{T1}{ptm}{m}{n}
\rput(10.527812,0.82968754){$D$}
\usefont{T1}{ptm}{m}{n}
\rput(10.371095,3.6896875){$E$}
\usefont{T1}{ptm}{m}{n}
\rput(3.8720312,1.0896875){$34^\circ$}
\usefont{T1}{ptm}{m}{n}
\rput(3.6520312,0.44968748){$62^\circ$}
\psarc[linewidth=0.024,linecolor=color194,arrowsize=0.05291667cm 2.0,arrowlength=1.4,arrowinset=0.4]{->}(4.07,1.0396874){0.49}{336.25052}{58.24052}
\psarc[linewidth=0.024,linecolor=color194,arrowsize=0.05291667cm 2.0,arrowlength=1.4,arrowinset=0.4]{<-}(3.56,0.80968744){0.94}{294.44397}{0.0}
\psframe[linewidth=0.04,linecolor=color194,dimen=outer,fillstyle=solid,fillcolor=color649b](10.38,-3.2103126)(0.0,-3.3903124)
\psline[linewidth=0.04](2.28,-2.886875)(2.6,-2.886875)(2.6,-3.226875)
\psline[linewidth=0.04](9.7,0.51312506)(9.7,1.153125)(10.05669,1.1505005)
\psline[linewidth=0.04cm](9.72,0.53312504)(10.0,0.53312504)
\psline[linewidth=0.04cm](9.72,0.83312505)(10.0,0.83312505)
\end{pspicture} 
}\\
Note: diagram not drawn to scale
\end{center}
}
{
\westep{To determine height $CE$, first calculate lengths $DE$ and $CD$}
$\triangle BDE$ and $\triangle BDC$ are both right-angled triangles. In each of the triangles, the length $BD$ is known. Therefore we can calculate the sides of the triangles.

\westep{Calculate $CD$}
The length $AC$ is given. $CABD$ is a rectangle so $BD = AC = 200$ m.\\
In $\triangle CBD$, 
\begin{eqnarray*}
tan~ C\hat{B}D &=& \frac{CD}{BD}\\
\therefore CD&=&BD\times tan~ C\hat{B}D \\
&=& 200\times tan~  62^{\circ} \\
&=& 376\mbox{ m}
\end{eqnarray*}

\westep{Calculate $DE$}
In $\triangle DBE$,
\begin{eqnarray*}
tan~  D\hat{B}E &=& \frac{DE}{BD}\\
\therefore DE&=&BD\times tan~ D\hat{B}E \\
&=& 200\times tan~  34^\circ \\
&=&135\mbox{ m}
\end{eqnarray*}

\westep{Add the two heights to get the final answer} 
The height of the tower $CE=CD+DE=135 \mbox{ m}+376\mbox{ m}=511\mbox{ m}$.
}
\end{wex}


\begin{wex}{Building plan}
{Mr Nkosi has a garage at his house and he decides to add a corrugated iron roof to the side of the garage. The garage is $4$ m high, and his sheet for the roof is $5$ m long. If the angle of the roof is $5^\circ$, how high must he build the wall $BD$? Give the answer correct to $1$ decimal place.
\\
\begin{center}
\scalebox{1} % Change this value to rescale the drawing.
{
\begin{pspicture}(0,-1.783125)(9.368906,1.783125)
\definecolor{color194}{rgb}{0.2,0.2,0.2}
\definecolor{color247b}{rgb}{0.8,0.8,0.8}
\psframe[linewidth=0.002,linecolor=white,linestyle=dotted,dotsep=0.16cm,dimen=outer,fillstyle=solid,fillcolor=color247b](8.471094,0.67968744)(8.311094,-1.6403126)
\psframe[linewidth=0.04,linecolor=white,dimen=outer,fillstyle=solid,fillcolor=color247b](4.3910937,0.25968745)(3.8510938,-1.6203126)
\psframe[linewidth=0.04,linecolor=white,dimen=outer,fillstyle=solid,fillcolor=color247b](1.2910937,0.25968745)(0.75109375,-1.6603125)
\psframe[linewidth=0.002,linecolor=white,linestyle=dotted,dotsep=0.16cm,dimen=outer,fillstyle=solid,fillcolor=color247b](4.3910937,1.3196875)(0.79109377,0.19968745)
% \usefont{T1}{ptm}{m}{n}
\rput(2.5167189,-0.20031255){Garage}
% \usefont{T1}{ptm}{m}{n}
\rput(7.78125,1.2396874){Roof}
% \usefont{T1}{ptm}{m}{n}
\rput(8.949843,-0.50031257){Wall}
\psline[linewidth=0.04cm](0.7710937,1.2996874)(0.7710937,-1.6003126)
\psline[linewidth=0.04cm](0.75109375,1.2996874)(4.351094,1.2996874)
\psline[linewidth=0.04cm](4.3710938,1.2996874)(4.3710938,-1.6003126)
\psline[linewidth=0.04cm](0.7710937,-1.6003126)(1.2710937,-1.6003126)
\psline[linewidth=0.04cm](1.2710937,-1.6003126)(1.2710937,0.19968745)
\psline[linewidth=0.04cm](1.2710937,0.19968745)(3.8710938,0.19968745)
\psline[linewidth=0.04cm](3.8710938,0.19968745)(3.8710938,-1.6003126)
\psline[linewidth=0.04cm](3.8710938,-1.6003126)(4.3710938,-1.6003126)
\psline[linewidth=0.024cm,linecolor=color194](4.3710938,1.2996874)(8.44,0.663125)
\psline[linewidth=0.027999999cm,linecolor=color194,linestyle=dashed,dash=0.16cm 0.16cm](8.4,0.663125)(4.2510934,0.6996874)
\psline[linewidth=0.04cm,linecolor=color194](8.291094,0.67968744)(8.291094,-1.6203126)
\psline[linewidth=0.04cm,linecolor=color194](8.451094,0.67968744)(8.451094,-1.6203126)
\psline[linewidth=0.018cm,linecolor=color194](8.46,-1.616875)(0.9710938,-1.6003126)
% \usefont{T1}{ptm}{m}{n}
\rput(0.3,0.07968745){$4$ m}
% \usefont{T1}{ptm}{m}{n}
\rput(6.3395314,1.2596875){$5$ m}
% \usefont{T1}{ptm}{m}{n}
\rput(8.906875,0.6996874){$ B$}
% \usefont{T1}{ptm}{m}{n}
\rput(4.6251564,0.41968745){$A$}
% \usefont{T1}{ptm}{m}{n}
\rput(6.073125,0.8796874){$5^\circ$}
% \usefont{T1}{ptm}{m}{n}
\rput(4.5270314,1.5796875){$C$}
\psdots[dotsize=0.12,linecolor=color194](4.3710938,0.6996874)
\psdots[dotsize=0.12,linecolor=color194](4.3710938,1.2996874)
% \usefont{T1}{ptm}{m}{n}
\rput(8.858906,-1.5803125){$ D$}
\psdots[dotsize=0.1378129,linecolor=color194](8.271093,0.67968744)
\psline[linewidth=0.04](4.38,0.943125)(4.64,0.943125)(4.64,0.703125)(4.64,0.703125)(4.64,0.723125)
\end{pspicture} 
}
\end{center}

}{
\westep{Identify opposite and adjacent sides and hypotenuse}
$\triangle ABC$ is right-angled. The hypotenuse and an angle are known therefore we can calculate $AC$. The height of the wall $BD$ is then the height of the garage minus $AC$.
\begin{eqnarray*}
 sin~ A\hat{B}C &=& \frac{AC}{BC} \\
\therefore AC &=& BC \times sin~ A\hat{B}C\\
&=& 5 sin~ 5^{\circ}\\
&=& 0,4\mbox{ m}\\
\\
\therefore BD&=& 4\mbox{ m}-0,4\mbox{ m}\\
&=& 3,6\mbox{ m}
\end{eqnarray*}


\westep{Write the final answer}  
Mr Nkosi must build his wall to be $3,6$ m high.
}
\end{wex}

\begin{exercises}{}
{
\begin{enumerate}[noitemsep, label=\textbf{\arabic*}. ] 

\item A boy flying a kite is standing $30~$m from a point directly under the kite. If the kite's string is $50~$m long, find the angle of elevation of the kite.
\item What is the angle of elevation of the sun when a tree $7,15$ m tall casts a shadow $10,1$ m long?
\item From a distance of $300$ m, Susan looks up at the top of a lighthouse. The angle of elevation is $5^{\circ}$. Determine the height of the lighthouse to the nearest metre.
\item A ladder of length $25$ m is resting against a wall, the ladder makes an angle $37^{\circ}$ to the wall. Find the distance between the wall and the base of the ladder. 

\end{enumerate}

\practiceinfo
\par 
\par \begin{tabular}[h]{cccccc}
(1.) 00kt&  (2.) 00ku& (3.) 00kv& (4.) 00ms \end{tabular}
    
}
\end{exercises} 

\section{Defining ratios in the Cartesian plane}

We have defined the trigonometric functions using right-angled
triangles. We can extend these definitions to any angle, noting that
the definitions do not rely on the lengths of the sides of the
triangle, but on the size of the angle only. So if we plot any point
on the Cartesian plane and then draw a line from the origin to that
point, we can work out the angle of that line. In the figure below
points $P$ and $Q$ have been plotted. A line from the origin ($O$) to
each point is drawn. The dotted lines show how we can construct
right-angled triangles for each point. The dotted line must always be
drawn to the $x$-axis. Now we can find the angles $A$ and $B$. We can
also extend the definitions of the reciprocals in the same way:


\setcounter{subfigure}{0}
\begin{center}

\scalebox{1} % Change this value to rescale the drawing.
{
\begin{pspicture}(0,-4.116719)(8.559063,4.1567187)
\definecolor{color4644b}{rgb}{0.7568627450980392,0.7568627450980392,0.7568627450980392}
\rput{127.209015}(6.6994696,-3.0593522){\psframe[linewidth=0.04,linecolor=color4644b,dimen=outer,fillstyle=solid,fillcolor=color4644b](4.2116795,0.34230337)(4.0061655,-0.07666689)}
\rput{0.05210805}(0.0,-0.0037759799){\psframe[linewidth=0.04,linecolor=color4644b,dimen=outer,fillstyle=solid,fillcolor=color4644b](4.2998266,0.5166024)(4.00418,-0.0911309)}
\rput{0.05210805}(0.0,-0.0039762384){\psframe[linewidth=0.04,linecolor=color4644b,dimen=outer,fillstyle=solid,fillcolor=color4644b](4.460063,0.3966389)(4.284289,-0.07118539)}
\rput{87.87539}(4.296069,-4.472766){\psframe[linewidth=0.04,linecolor=color4644b,dimen=outer,fillstyle=solid,fillcolor=color4644b](4.571679,0.20230302)(4.3661656,-0.21666722)}
\rput{32.417133}(0.7902694,-1.9834906){\psframe[linewidth=0.04,linecolor=color4644b,dimen=outer,fillstyle=solid,fillcolor=color4644b](3.9364834,0.65755856)(3.677197,0.07755861)}
\rput{-56.56522}(1.8740156,3.179452){\psarc[linewidth=0.04,fillstyle=solid,fillcolor=color4644b](3.8915973,-0.15175243){0.831297}{60.043613}{171.69456}}
\rput(4.0,-0.11671885){\psaxes[linewidth=0.04,arrowsize=0.05291667cm 2.0,arrowlength=1.4,arrowinset=0.4,ticksize=0.15cm]{<->}(0,0)(-4,-4)(4,4)}
\psline[linewidth=0.04cm,dotsize=0.07055555cm 2.0]{-*}(4.04,-0.09671885)(6.0,2.9167187)
\psline[linewidth=0.04cm,dotsize=0.07055555cm 2.0]{-*}(4.02,-0.14328125)(2.0,2.8832812)
\psline[linewidth=0.04cm,linestyle=dashed,dash=0.16cm 0.16cm](5.98,2.9167187)(6.0,-0.08328125)
\psline[linewidth=0.04cm,linestyle=dashed,dash=0.16cm 0.16cm](2.0,2.8364403)(2.0,-0.09671885)
\psline[linewidth=0.04,fillstyle=solid](2.0,0.17671885)(2.26,0.17671885)(2.26,-0.07568518)(2.26,-0.09671885)
\psline[linewidth=0.04,fillstyle=solid](5.72,-0.08328115)(5.72,0.17518039)(5.98,0.17518039)(5.98,0.19671886)
\rput{-279.77182}(3.025072,-3.2106922){\psarc[linewidth=0.04](3.4179945,0.18995206){0.66433096}{15.686174}{126.06392}}
\rput{-86.49843}(3.894729,4.3688884){\psarc[linewidth=0.04](4.269561,0.11427801){0.9286654}{72.6542}{148.98633}}
\usefont{T1}{ptm}{m}{n}
\rput(4.849375,0.31328115){$A$}
\usefont{T1}{ptm}{m}{n}
\rput(3.8090625,-0.31828135){$0$}
\usefont{T1}{ptm}{m}{n}
\rput(3.2445314,0.25328115){$x$}
\usefont{T1}{ptm}{m}{n}
\rput(6.606094,3.1332812){$P(2;3)$}
\usefont{T1}{ptm}{m}{n}
\rput(1.3735937,3.1532812){$Q(-2;3)$}
\usefont{T1}{ptm}{m}{n}
\rput(8.184531,0.09328115){$x$}
\usefont{T1}{ptm}{m}{n}
\rput(4.2248435,3.9532814){$y$}
\usefont{T1}{ptm}{m}{n}
\rput(4.1221876,0.39328116){$B$}
\end{pspicture} 
}
\end{center}
From the coordinates of $P(2;3)$, we know the length of the side opposite $\hat{A}$ is $3$ and the length of the adjacent side is $2$. Using $tan~ \hat{A}=\frac{\mbox{\footnotesize opposite}}{\mbox{\footnotesize adjacent}} = \frac{3}{2}$, we calculate that $\hat{A}=56,3^{\circ}$.\par

We can also use the Theorem of Pythagoras to calculate the hypotenuse of the triangle and then calculate $\hat{A}$ using $sin~ \hat{A} = \frac{\mbox{\footnotesize opposite}}{\mbox{\footnotesize hypotenuse}}$ or $cos~ \hat{A} = \frac{\mbox{\footnotesize adjacent}}{\mbox{\footnotesize hypotenuse}}$. \par
Consider point $Q(-2;3)$. We define $\hat{B}$ as the angle formed between the line $OQ$ and the positive $x$-axis. This is called the standard position of an angle. Angles are always measured from the positive $x$-axis in an anti-clockwise direction. Let $\hat{x}$ be the angle formed between the line $OQ$ and the negative $x$-axis such that $\hat{B} + \hat{x} = 180^{\circ}$.
\par

From the coordinates of $Q(-2;3)$, we know the length of the side opposite $\hat{x}$ is $3$ and the length of the adjacent side is $2$. Using $tan~ \hat{x}=\frac{\mbox{\footnotesize opposite}}{\mbox{\footnotesize adjacent}} = \frac{3}{2}$ we calculate that $\hat{x}=56,3^{\circ}$.
\\Therefore $\hat{B}=180^{\circ} - \hat{x} = 123,7^{\circ}$.\par

Similarly, an alternative method is to calculate the hypotenuse using the Theorem of Pythagoras and calculate $\hat{x}$ using $sin~ \hat{x} = \frac{\mbox{\footnotesize opposite}}{\mbox{\footnotesize hypotenuse}}$ or $cos~ \hat{x} = \frac{\mbox{\footnotesize adjacent}}{\mbox{\footnotesize hypotenuse}}$. \par
If we were to draw a circle centred on the origin $(O)$ and passing through point $P$, then the length from the origin to point $P$ is the radius of the circle, which we denote $r$. We can rewrite all the trigonometric functions in terms of $x$, $y$ and $r$.


The general definitions for the trigonometric functions are:

\begin{equation*}
\begin{array}{cccccc}\hfill sin~ \theta & =& \frac{y}{r}\hfill & cosec~ \theta & =& \frac{r}{y}\hfill \\
 \hfill cos~ \theta & =& \frac{x}{r}\hfill & sec~ \theta & =& \frac{r}{x}\hfill \\
 \hfill tan~ \theta & =& \frac{y}{x}\hfill & cot~ \theta & =& \frac{x}{y}\hfill \end{array}
\end{equation*}


The Cartesian plane is divided into $4$ quadrants in an anti-clockwise direction as shown in the diagram below. Notice that $r$ is always positive but the values of $x$ and $y$ change depending on the position of the point in the Cartesian plane. As a result, the trigonometric ratios can be positive or negative. The letters $C$, $A$, $S$ and $T$ indicate which of the functions are positive in each quadrant: \\
\\
Quadrant I: All ratios are positive.\\
Quadrant II: $y$ values are positive therefore sine and cosine are positive.\\
Quadrant III: Both $x$ and $y$ values are negative therefore tangent and cotangent are positive. \\
Quadrant IV: $x$ values are positive therefore cosecant and secant are positive.\par

\textbf{Note:} the hypotenuse  $r$ is a length, and is therefore always positive. 
% \Note{$r$ is a length and therefore is always positive.}
\begin{center}
\scalebox{1} % Change this value to rescale the drawing.
{
\begin{pspicture}(0,-4.1584377)(8.42125,4.1984377)
\rput(4.0,-0.1584375){\psaxes[linewidth=0.04,arrowsize=0.05291667cm 2.0,arrowlength=1.4,arrowinset=0.4,labels=none,ticks=none,ticksize=0.10583333cm]{<->}(0,0)(-4,-4)(4,4)}
% \usefont{T1}{ptm}{m}{n}
\rput(6.0853124,2.2715626){Quadrant I}
% \usefont{T1}{ptm}{m}{n}
\rput(1.7453125,2.3115625){Quadrant II}
% \usefont{T1}{ptm}{m}{n}
\rput(1.8053125,-2.3084376){Quadrant III}
% \usefont{T1}{ptm}{m}{n}
\rput(6.3515625,-2.2684374){Quadrant IV}
% \usefont{T1}{ptm}{m}{n}
\rput(4.7403126,0.8315625){A }
% \usefont{T1}{ptm}{m}{n}
\rput(4.7723436,0.4315625){all}
% \usefont{T1}{ptm}{m}{n}
\rput(3.2923439,0.8515625){S}
% \usefont{T1}{ptm}{m}{n}
\rput(3.2390625,0.4715625){$sin~\theta$}
% \usefont{T1}{ptm}{m}{n}
\rput(3.2109375,-0.8484375){T}
% \usefont{T1}{ptm}{m}{n}
\rput(3.1639063,-1.1684375){$tan~\theta$}
% \usefont{T1}{ptm}{m}{n}
\rput(4.7857814,-0.8484375){C}
% \usefont{T1}{ptm}{m}{n}
\rput(4.815,-1.1484375){$cos~\theta$}
\usefont{T1}{ptm}{m}{n}
\rput(8.275469,-0.1084375){$x$}
\usefont{T1}{ptm}{m}{n}
\rput(4.1557813,4.0715623){$y$}
\usefont{T1}{ptm}{m}{n}
\rput(7.810625,0.0715625){$0^{\circ}$}
\usefont{T1}{ptm}{m}{n}
\rput(7.60875,-0.4484375){$360^{\circ}$}
\usefont{T1}{ptm}{m}{n}
\rput(4.3410935,3.5515625){$90^{\circ}$}
\usefont{T1}{ptm}{m}{n}
\rput(0.32265624,0.0715625){$180^{\circ}$}
\usefont{T1}{ptm}{m}{n}
\rput(3.529375,-3.9084375){$270^{\circ}$}

\usefont{T1}{ptm}{m}{n}
\rput(3.8,-0.4084375){$0$}
\end{pspicture} 
}
\end{center}

This diagram is known as the CAST diagram.
\par
\mindsetvid{Trigonometry on the Cartesian plane}{VMbhy}

\subsubsection{Special angles in the Cartesian plane}
When working in the Cartesian plane we include two other special
angles in right-angled triangles: $0^{\circ}$ and $90^{\circ}$. These
are special since we can usually not have an angle of $0^{\circ}$ or
another angle (in addition to the right angle) of $90^{\circ}$.

Notice that when $\theta = 0^{\circ}$, the length of the opposite side
is equal to $0$, therefore
\begin{equation*}
  sin~ 0^{\circ} =
  \dfrac{\mbox{\footnotesize opposite}}{\mbox{\footnotesize hypotenuse}} =
  \dfrac{0}{\mbox{\footnotesize hypotenuse}} =
  0.
\end{equation*}

When $\theta = 90^{\circ}$, the length of the adjacent side is equal
to $0$, therefore
\begin{equation*}
  cos~ 90^{\circ} =
  \dfrac{\mbox{\footnotesize adjacent}}{\mbox{\footnotesize hypotenuse}} =
  \dfrac{0}{\mbox{\footnotesize hypotenuse}} =
  0.
\end{equation*}

Using the definition $tan~ \theta=\dfrac{sin~ \theta}{cos~ \theta}$ we
see that for $\theta =  0^{\circ}$, $sin~ 0^{\circ}=0$, therefore
\begin{equation*}
  tan~ 0^{\circ} =
  \dfrac{0}{cos~ 0^{\circ}} =
  0.
\end{equation*}
For $\theta =  90^{\circ}$, $cos~ 90^{\circ}=0$, therefore
\begin{equation*}
  tan~ 90^{\circ} =
  \dfrac{sin~ 90^{\circ}}{0} =
  \mbox{undefined.}
\end{equation*}

\begin{table}[H]
% \begin{table}[H]
% \\ '' '0'
\begin{center}
\setlength{\extrarowheight}{2.5pt}

\begin{tabular}{|c|c|c|c|c|c|}\hline
$\theta$
&
${0}^{\circ }$
&
${30}^{\circ }$
&
${45}^{\circ }$
&
${60}^{\circ }$
&
${90}^{\circ }$



\\ \hline
%--------------------------------------------------------------------
$cos~ \theta $
 &
$1$
&
$\frac{\sqrt{3}}{2}$
&
$\frac{1}{\sqrt{2}}$
&
$\frac{1}{2}$
&
$0$


\\ \hline
%--------------------------------------------------------------------
$sin~ \theta $
&
$0$
&
$\frac{1}{2} $
&
$\frac{1}{\sqrt{2}}$
&
$\frac{\sqrt{3}}{2}$
&
$1$

%  make-rowspan-placeholders
\\ \hline
%--------------------------------------------------------------------
$tan~ \theta $
&
$0$
&
$\frac{1}{\sqrt{3}}$
&
$1$ &
$\sqrt{3}$
&
undef


% make-rowspan-placeholders
\\ \hline
%--------------------------------------------------------------------
\end{tabular}
\end{center}
% \begin{center}{\small\bfseries Table 14.3}\end{center}
% \begin{caption}{\small\bfseries Table 14.3}\end{caption}
\end{table}



\begin{wex}{Ratios in the Cartesian plane}
{
\begin{minipage}{\textwidth}
$P(-3;4)$ is a point in the Cartesian plane. $X\hat{O}P=\theta$,
  where $X$ is a point on the $x$-axis. Without using a calculator,
  determine the value of:
\begin{enumerate}[noitemsep, label=\textbf{\arabic*}. ] 
\item $cos~ \theta$
\item $3 tan~ \theta$
\item $\frac{1}{2} cosec~ \theta$
\end{enumerate}
\end{minipage}
}
{
\westep{Sketch point $P$ in the Cartesian plane and label $\theta$} 
\begin{center}
\scalebox{1} % Change this value to rescale die drawing.
{
\begin{pspicture}(0,-3.1584375)(6.44125,3.1984375)
\rput(3.0,-0.1584375){\psaxes[linewidth=0.04,arrowsize=0.05291667cm 2.0,arrowlength=1.4,arrowinset=0.4,labels=none,ticks=none,ticksize=0.10583333cm]{<->}(0,0)(-3,-3)(3,3)}
\usefont{T1}{ptm}{m}{n}
\rput(6.295469,0.0115625){$x$}
\usefont{T1}{ptm}{m}{n}
\rput(3.1757812,3.0715625){$y$}
\psline[linewidth=0.04cm,dotsize=0.07055555cm 2.0]{-*}(3.02,-0.1784375)(1.6,1.8615625)
\usefont{T1}{ptm}{m}{n}
\rput(1.3670312,2.2715626){$P(-3;4)$}
\rput{-45.048485}(0.95434314,2.1533978){\psarc[linewidth=0.04]{->}(3.073445,-0.073917985){0.6961048}{36.728504}{180.0}}
\usefont{T1}{ptm}{m}{n}
\rput(3.150625,-0.4484375){$0$}
\usefont{T1}{ptm}{m}{n}
\rput(3.2709374,0.1915625){$\theta$}
\end{pspicture} 
}
\end{center}

\westep{Use Pythagoras to calculate $r$}

\begin{eqnarray*}
  r^{2} &=& x^{2} + y^{2} \\
        &=& (-3)^{2} + (4)^{2} \\
        &=& 25 \\
  \therefore r &=& 5
\end{eqnarray*}
\westep{Substitute values for $x$, $y$ and $r$ into the required ratios}
\vspace{1ex}
\begin{minipage}{\textwidth}
\begin{enumerate}[itemsep=5pt, label=\textbf{\arabic*}. ] 
\item $cos~ \theta = \dfrac{x}{r} = -\dfrac{3}{5}$
\item $3 tan~ \theta = 3\left(\dfrac{y}{x}\right) = 3\left(\dfrac{4}{-3}\right) = -4 $
\item $\dfrac{1}{2} cosec~ \theta = \dfrac{1}{2}\left(\dfrac{r}{y}\right) = \dfrac{1}{2}\left(\dfrac{5}{4}\right) = \dfrac{5}{8} $
\end{enumerate}
\end{minipage}
}
\end{wex}


\begin{wex}{Ratios in the Cartesian plane}
{
\begin{minipage}{\textwidth}
$X\hat{O}K = \theta$ is an angle in the third quadrant and $K$ is the point $(-5;y)$. $OK$ is $13$ units. Determine without using a calculator:
  \begin{enumerate}[noitemsep, label=\textbf{\arabic*}. ] 
   \item The value of $y$
\item Prove that $tan~^{2}~\theta + 1 = sec~^{2}~\theta$
  \end{enumerate}
\end{minipage}
}
{
\westep{Sketch point $K$ in the Cartesian plane and label $\theta$} 
\begin{center}
\scalebox{1} % Change this value to rescale die drawing.
{
\begin{pspicture}(0,-3.1584375)(6.44125,3.1984375)
\rput(3.0,-0.1584375){\psaxes[linewidth=0.04,arrowsize=0.05291667cm 2.0,arrowlength=1.4,arrowinset=0.4,labels=none,ticks=none,ticksize=0.10583333cm]{<->}(0,0)(-3,-3)(3,3)}
\usefont{T1}{ptm}{m}{n}
\rput(6.295469,0.0115625){$x$}
\usefont{T1}{ptm}{m}{n}
\rput(3.1757812,3.0715625){$y$}
\psline[linewidth=0.04cm,dotsize=0.07055555cm 2.0]{-*}(2.963227,-0.18549357)(1.6,-2.1384375)
\usefont{T1}{ptm}{m}{n}
\rput(1.2170312,-2.5084374){$K(-5;y)$}
\rput{-45.048485}(1.0897402,2.0403726){\psarc[linewidth=0.04]{->}(3.0048735,-0.2936738){0.7345531}{56.389397}{272.69913}}
\usefont{T1}{ptm}{m}{n}
\rput(3.150625,-0.4484375){$0$}
\usefont{T1}{ptm}{m}{n}
\rput(2.4109375,0.5115625){$\theta$}
\usefont{T1}{ptm}{m}{n}
\rput(1.9303125,-1.1484375){$13$}
\end{pspicture} 
}
\end{center}
\westep{Use Pythagoras to calculate $y$}
\begin{equation*}
 \begin{array}{ccl}
    \hfill r^{2} &= & x^{2} + y^{2} \hfill \\
\hfill y^{2} &=& r^{2} - x^{2}\hfill \\
\hfill  &=& (13)^{2} - (-5)^{2} \hfill \\
\hfill  &=& 169-25 \hfill \\
\hfill  &=& 144 \hfill \\
\hfill  y&=& \pm 12\hfill 

\end{array}
\end{equation*}
Given that $\theta$ lies in the third quadrant, $y$ must be negative.\\
\begin{equation*}
 \therefore y = -12
\end{equation*}

\westep{Substitute values for $x$, $y$ and $r$ and simplify}
\begin{equation*}
\begin{array}{rclrcl}
 \hfill \mbox{LHS} &=&  tan~^{2}~\theta + 1 	&	\mbox{RHS}	&=&  	sec~^{2}~\theta \hfill \\
\hfill &=& \left(\dfrac{y}{x}\right)^{2} + 1  	&	&=& \left(\dfrac{r}{x}\right)^{2} \hfill \vspace{5pt} \\
\hfill&=&  \left(\dfrac{-12}{-5}\right)^{2} + 1 	&	&=& \left(\dfrac{13}{-5}\right)^{2} \hfill \vspace{5pt}\\
\hfill&=&  \left (\dfrac{144}{25}\right) + 1  	&	&=& \dfrac{169}{25} \hfill \vspace{5pt}\\
\hfill &=& \dfrac{144 + 25}{25} 		&	& & \hfill \vspace{5pt}\\
\hfill &=& \dfrac{169}{25}  			&	& & \hfill \vspace{5pt}\\
\hfill \therefore \mbox{LHS} &=& \mbox{RHS}&& \hfill
\end{array}
\end{equation*}
}
\end{wex}


\begin{exercises}{}
{
\textbf{Important:} whenever you have to solve trigonometric problems without a calculator, always make a sketch.\\

  \begin{enumerate}[itemsep=5pt, label=\textbf{\arabic*}. ]
   \item $B$ is a point in the Cartesian plane. Determine without using a calculator:

% \begin{center}
% \scalebox{1} % Change this value to rescale the drawing.
% {
% \begin{pspicture}(0,-3.1584375)(6.44125,3.1984375)
% \rput(3.0,-0.1584375){\psaxes[linewidth=0.04,arrowsize=0.05291667cm 2.0,arrowlength=1.4,arrowinset=0.4,labels=none,ticks=none,ticksize=0.10583333cm]{<->}(0,0)(-3,-3)(3,3)}
% \usefont{T1}{ppl}{m}{n}
% \rput(6.295469,0.0115625){$x$}
% \usefont{T1}{ppl}{m}{n}
% \rput(3.1757812,3.0715625){$y$}
% \rput(3.4, -0.4){$O$}
% \psline[linewidth=0.04cm,fillcolor=black,dotsize=0.07055555cm 2.0]{-*}(3.003227,-0.16549358)(3.78,-1.8384376)
% \usefont{T1}{ppl}{m}{n}
% \rput(4.577031,-1.9084375){$B(1;-3)$}
% \rput{-45.048485}(1.0590175,2.0862904){\psarc[linewidth=0.04](3.0448735,-0.2336738){0.7345531}{50.714928}{340.5804}}
% \usefont{T1}{ppl}{m}{n}
% \rput(2.393125,0.5115625){$\beta$}
% \end{pspicture} 
% }
% \end{center}
\begin{enumerate}[noitemsep, label=\textbf{(\alph*)} ]
\item $OB$
\item $cos~ \beta$
\item $cosec~ \beta$
\item $tan~ \beta$
\end{enumerate}


\scalebox{1} % Change this value to rescale the drawing.
{
\begin{pspicture}(0,-3.1567187)(6.5990624,3.1967187)
\rput(3.0,-0.15671875){\psaxes[linewidth=0.04,arrowsize=0.05291667cm 2.0,arrowlength=1.4,arrowinset=0.4,labels=none,ticks=none,ticksize=0.10583333cm]{<->}(0,0)(-2.5,-2.5)(2.5,2.5)}
\usefont{T1}{ptm}{m}{n}
\rput(6.224531,0.01328125){$x$}
\usefont{T1}{ptm}{m}{n}
\rput(3.1245313,2.9932814){$y$}
\psline[linewidth=0.04cm](3.0,-0.17671876)(4.0,-2.2167187)
\psline[linewidth=0.04cm,linestyle=dashed,dash=0.16cm 0.16cm](3.98,-2.2167187)(4.02,-0.09671875)
\psdots[dotsize=0.12](4.0,-2.2367187)
\psarc[linewidth=0.04,arrowsize=0.05291667cm 2.0,arrowlength=1.4,arrowinset=0.4]{->}(3.02,-0.15671875){0.4}{0.0}{284.03625}
\psframe[linewidth=0.04,dimen=outer](4.04,-0.13671875)(3.76,-0.41671875)
\usefont{T1}{ptm}{m}{n}
\rput(3.9645312,0.1){$X$}
\usefont{T1}{ptm}{m}{n}
% \rput(2.9,-0.4){$0$}
\usefont{T1}{ptm}{m}{n}
\rput(2.5445313,0.33328125){$\beta$}
\usefont{T1}{ptm}{m}{n}
\rput(4.1445312,-2.5067186){$B(1;-3)$}
\end{pspicture} 
}


\item If $sin~ \theta= 0,4$ and $\theta$ is an obtuse angle, determine:
% \begin{center}
% \scalebox{1} % Change this value to rescale the drawing.
% {
% \begin{pspicture}(0,-3.1584375)(6.44125,3.1984375)
% \rput(3.0,-0.1584375){\psaxes[linewidth=0.04,arrowsize=0.05291667cm 2.0,arrowlength=1.4,arrowinset=0.4,labels=none,ticks=none,ticksize=0.10583333cm]{<->}(0,0)(-3,-3)(3,3)}
% \usefont{T1}{ppl}{m}{n}
% \rput(6.295469,0.0115625){$x$}
% \usefont{T1}{ppl}{m}{n}
% \rput(3.1757812,3.0715625){$y$}
% \psline[linewidth=0.04cm,fillcolor=black](3.0,-0.1784375)(0.8,0.7415625)
% \usefont{T1}{ppl}{m}{n}
% \rput(0.7292187,-0.4884375){$(x;2)$}
% \rput{-45.048485}(1.0440966,2.0498228){\psarc[linewidth=0.04](2.9934452,-0.23391798){0.6961048}{50.894093}{199.87982}}
% \usefont{T1}{ppl}{m}{n}
% \rput(3.150625,-0.4484375){$0$}
% \usefont{T1}{ppl}{m}{n}
% \rput(3.6509376,0.5915625){$\theta$}
% \psline[linewidth=0.04cm,linestyle=dashed,dash=0.16cm 0.16cm](0.8,0.7215625)(0.8,-0.1784375)
% \usefont{T1}{ptm}{m}{n}
% \rput(0.5234375,0.2315625){$2$}
% \end{pspicture} 
% }
% \end{center}
\begin{enumerate}[noitemsep, label=\textbf{(\alph*)} ]
% \setcounter{enumi}{4}
 \item $cos~ \theta$
\item $\sqrt{21} tan~ \theta$
\end{enumerate}
\item Given $tan~ \theta = \frac{t}{2}$, where $0^{\circ} \leq \theta \leq 90^{\circ}$.
Determine the following in terms of $t$:
\begin{enumerate}[noitemsep, label=\textbf{(\alph*)} ]
\item $sec~ \theta$
\item $cot~ \theta$
\item $cos^2~ \theta$
\item $tan^2~ \theta-sec^2~ \theta$
\end{enumerate}
\end{enumerate}
% Automatically inserted shortcodes - number to insert 2
\par \practiceinfo
\par \begin{tabular}[h]{cccccc}
% Question 1
(1.)	00kw	&
% Question 2
(2.)	00kx	& (3.) 0242	& 
\end{tabular}
% Automatically inserted shortcodes - number inserted 2
}
\end{exercises}

\summary{VMdwh}
\begin{itemize}[noitemsep]
\item We can define three trigonometric ratios for right-angled triangles: sine ($sin$), cosine ($cos$) and tangent ($tan$).
\item Each of these ratios have a reciprocal: cosecant ($cosec$), secant ($sec$) and cotangent ($cot$).
\item We can use the principles of solving equations and the trigonometric ratios to help us solve simple trigonometric equations.
\item We can solve problems in two dimensions that involve right-angled triangles.
\item For some special angles, we can easily find the values of $sin$, $cos$ and $tan$ without using a calculator.
\item We can extend the definitions of the trigonometric functions to any angle.
\item Trigonometry is used to help us solve problems in two dimensions, such as finding the height of a building.
\end{itemize}


\begin{eocexercises}{}

\begin{enumerate}[itemsep=6pt, label=\textbf{\arabic*}. ] 
\item Without using a calculator determine the value of 
\begin{equation*}
sin~ 60^{\circ} cos~ 30^{\circ}-cos~ 60^{\circ}sin~ 30^{\circ} + tan~ 45^{\circ}
\end{equation*}
\item If $3 tan~ \alpha = -5$ and $0^{\circ} < \alpha < 270^{\circ}$, use a sketch to determine:
    \begin{enumerate}[noitemsep, label=\textbf{(\alph*)} ]
    \item $cos~ \alpha$
    \item $tan~^{2}~\alpha - sec~^{2}~\alpha$
    \end{enumerate}
\item Solve for $\theta$ if $\theta$ is a positive, acute angle:
    \begin{enumerate}[noitemsep, label=\textbf{(\alph*)} ]
    \item $2 sin~ \theta = 1,34$
    \item $1 - tan~ \theta = -1$
    \item $cos~ 2\theta = sin~ 40^{\circ}$ 
    \item $\dfrac{sin~ \theta}{cos~ \theta}= 1$
    \end{enumerate}


\item Calculate the unknown lengths in the diagrams below:
\begin{center}
\scalebox{1}  
{ 
\begin{pspicture}(0,-2.0390613)(8.035,2.0253136) 
\psline[linewidth=0.04cm](0.02,0.02093862)(3.06,0.02093862) 
\psline[linewidth=0.04cm](0.02,0.02093862)(0.02,-2.0190613) 
\psline[linewidth=0.04cm](0.04,-1.9990613)(3.04,0.02093862) 
\psframe[linewidth=0.04,dimen=outer](0.26,0.04093862)(0.0,-0.21906137) 
\psline[linewidth=0.04cm](0.9944844,1.4531701)(3.04,0.04093862) 
\psline[linewidth=0.04cm](0.96,1.4209386)(0.024381146,0.02542873) 
\rput{-33.90198}(-0.5512336,0.79249614){\psframe[linewidth=0.04,dimen=outer](1.1544285,1.4305158)(0.8944285,1.1705158)} 
\psline[linewidth=0.04cm](2.2384837,1.992771)(3.04,0.04093862) 
\psline[linewidth=0.04cm](2.238,1.9929386)(1.0103352,1.4765087) 
\rput{-66.64198}(-0.3803531,3.117777){\psframe[linewidth=0.04,dimen=outer](2.3111107,1.9781737)(2.0511107,1.7181737)} 
% \usefont{T1}{ptm}{m}{n} 
\rput(2.2601562,-0.21406138){$30^{\circ}$} 
% \usefont{T1}{ptm}{m}{n} 
\rput(2.3034375,0.22593862){$25^{\circ}$} 
% \usefont{T1}{ptm}{m}{n} 
\rput(2.4634376,0.77){$20^{\circ}$} 
% \usefont{T1}{ptm}{m}{n} 
\rput(1.8576562,-1.2540613){$16$ cm} 
% \usefont{T1}{ptm}{m}{it} 
\rput(1.3398438,0.20593862){$a$} 
% \usefont{T1}{ptm}{m}{it} 
\rput(1.9,1.1){$b$} 
% \usefont{T1}{ptm}{m}{it} 
\rput(1.58,1.9059386){$c$} 
\psline[linewidth=0.04cm](6.689009,-1.6241995)(6.0199666,-1.4501117) 
\rput{-90.46057}(4.6635084,7.686732){\psframe[linewidth=0.04,dimen=outer](6.2543488,1.6402806)(6.0343485,1.4202806)} 
\rput{-90.46057}(4.4806204,7.508202){\psframe[linewidth=0.04,dimen=outer](6.0743546,1.6417276)(5.8543544,1.4217275)} 
\rput{-104.9971}(9.7950115,4.4890895){\psframe[linewidth=0.04,dimen=outer](6.7298956,-1.4036404)(6.509896,-1.6236403)} 
\psline[linewidth=0.04cm](7.375193,1.6303898)(4.2549725,1.6154709) 
\psline[linewidth=0.04cm](6.055236,1.6410005)(6.030317,-1.4588993) 
\psline[linewidth=0.04cm](6.0304775,-1.4389)(4.2748113,1.5953108) 
\psline[linewidth=0.04cm](7.375193,1.6303898)(6.030317,-1.4588993) 
\psline[linewidth=0.04cm](7.375193,1.6303898)(6.689009,-1.6241995) 
% \usefont{T1}{ptm}{m}{it} 
\rput(4.880625,0.005938619){$d$} 
% \usefont{T1}{ptm}{m}{it} 
\rput(6.7414064,1.8259386){$e$} 
% \usefont{T1}{ptm}{m}{n} 
\rput(5.2226562,1.8259386){$5$ m} 
% \usefont{T1}{ptm}{m}{it} 
\rput(6.2709374,-1.7940614){$f$} 
% \usefont{T1}{ptm}{m}{it} 
\rput(7.2148438,-0.07406138){$g$} 
% \usefont{T1}{ptm}{m}{n} 
\rput(4.756875,1.3859386){ $50^{\circ}$} 
% \usefont{T1}{ptm}{m}{n} 
\rput(6.9460936,1.4059386){$60^{\circ}$} 
% \usefont{T1}{ptm}{m}{n} 
\rput(6.45,-1.2740613){$80^{\circ}$} 
\end{pspicture} 
}
\end{center}

\item In $\triangle PQR$, $PR=20$ cm, $QR=22$ cm and $P\hat{R}Q = 30^{\circ}$. The perpendicular line from $P$ to $QR$ intersects $QR$ at $X$. Calculate 
\begin{enumerate}[noitemsep, label=\textbf{(\alph*)} ]
\item the length $XR$ 
\item the length $PX$
\item the angle $Q\hat{P}X$ 
\end{enumerate} 
\item A ladder of length $15$ m is resting against a wall, the base of the ladder is $5$ m from the wall. Find the angle between the wall and the ladder. 
\item In the following triangle find the angle $A\hat{B}C$:
\begin{center}
\begin{pspicture}(0,-2.4701562)(5.49875,2.4701562) 
\pspolygon[linewidth=0.04](0.1665625,-1.7301563)(3.3665626,1.9698437)(5.1665626,-1.7301563)(4.1665626,-1.7301563) 
\psline[linewidth=0.04cm](3.3665626,1.9698437)(3.3665626,-1.7301563) 
\rput(3.3871875,2.2798438){$A$} 
\rput(5.3459377,-2.0201561){$B$} 
\rput(3.371875,-2.0201561){$C$} 
\rput(0.07546875,-2.0201561){$D$} 
\rput(3.6,0){$9$} 
\rput(2.7525,-2.3201563){$17$} 
\psline[linewidth=0.04cm,arrowsize=0.05291667cm 2.0,arrowlength=1.4,arrowinset=0.4]{->}(3.0665624,-2.3301563)(5.2665625,-2.3301563) 
\psline[linewidth=0.04cm,arrowsize=0.05291667cm 2.0,arrowlength=1.4,arrowinset=0.4]{->}(2.4665625,-2.3301563)(0.0665625,-2.3301563) 
\psline[linewidth=0.04cm](3.3665626,-1.5301563)(3.5665624,-1.5301563) 
\psline[linewidth=0.04cm](3.5665624,-1.5301563)(3.5665624,-1.7301563) 
\rput(0.8,-1.48){$41^{\circ}$} 
\end{pspicture} 
\end{center}
\item In the following triangle find the length of side $CD$:
\begin{center}
\begin{pspicture}(0,-2.2234375)(6.091875,2.2234375) 
\pspolygon[linewidth=0.04](0.1665625,-1.776875)(5.1665626,-1.776875)(5.1665626,1.823125) 
\psline[linewidth=0.04cm](3.4665625,-1.776875)(5.1665626,1.823125) 
\rput(5.2871876,2.033125){$A$} 
\rput(5.3459377,-2.066875){$B$} 
\rput(3.471875,-2.066875){$C$} 
\rput(0.07546875,-2.066875){$D$} 
\rput(5.4,0){$9$} 
\rput(4.490156,0.95){$15^{\circ}$} 
\rput(3.960156,-1.5){$35^{\circ}$} 
\psline[linewidth=0.04cm](4.9665626,-1.576875)(5.1665626,-1.576875) 
\psline[linewidth=0.04cm](4.9665626,-1.576875)(4.9665626,-1.776875) 
\end{pspicture}
\end{center} 
\item Given $A(5;0)$ and $B(11;4)$, find the angle between the line through $A$ and $B$ and the $x$-axis. 
\item Given $C(0;-13)$ and $D(-12;14)$, find the angle between the line through $C$ and $D$ and the $y$-axis. 


\item A right-angled triangle has hypotenuse $13$ mm. Find the length of the other two sides if one of the angles of the triangle is $50^{\circ}$.
\item One of the angles of a rhombus with perimeter $20$ cm is $30^{\circ}$. 
\begin{enumerate}[noitemsep, label=\textbf{(\alph*)} ]
\item Find the sides of the rhombus. 
\item Find the length of both diagonals. 
\end{enumerate} 
\item Captain Jack is sailing towards a cliff with a height of $10$ m. 
\begin{enumerate}[noitemsep, label=\textbf{(\alph*)} ] 
\item The distance from the boat to the top of the cliff is $30$ m. Calculate the angle of elevation from the boat to the top of the cliff (correct to the nearest integer).
\item If the boat sails $7$ m closer to the cliff, what is the new angle of elevation from the boat to the top of the cliff? 
\end{enumerate} 
\item Given the points, $E(5;0)$, $F(6;2)$ and $G(8;-2)$, find the angle $F\hat{E}G$. 
\item  A triangle with angles $40^{\circ}$, $40^{\circ}$ and $100^{\circ}$ has a perimeter of $20$ cm. Find the length of each side of the triangle. 

\end{enumerate}
% Automatically inserted shortcodes - number to insert 15
\par \practiceinfo
\par \begin{tabular}[h]{cccccc}
% Question 1
(1.)	00ky	&
% Question 2
(2.)	00kz	&
% Question 3
(3.)	00m0	&
% Question 4
(4.)	00m1	&
% Question 5
(5.)	00m2	&
% Question 6
(6.)	00m3	\\ % End row of shortcodes
% Question 7
(7.)	00m4	&
% Question 8
(8.)	00m5	&
% Question 9
(9.)	00m6	&
% Question 10
(10.)	00m7	&
% Question 11
(11.)	00m8	&
% Question 12
(12.)	00m9	\\ % End row of shortcodes
% Question 13
(13.)	00ma	&
% Question 14
(14.)	00mb	&
% Question 15
(15.)	00mc	&
\end{tabular}
% Automatically inserted shortcodes - number inserted 15
\end{eocexercises}
